% Performance Before Outcomes -- final edition, November 2026
% Single file: run pdflatex twice. Needs only the figures/ folder next to this file.
% Every data-derived number below is a macro generated by code/make_numbers.py from results/*.json.
\documentclass[9pt,twocolumn,letterpaper]{extarticle}
\usepackage[T1]{fontenc}
\usepackage[utf8]{inputenc}
\usepackage{mathpazo}
\linespread{1.04}
\usepackage[scaled=0.90]{helvet}
\usepackage{amsmath,amssymb,amsthm,mathtools,bm}
\usepackage[letterpaper,left=0.62in,right=0.62in,top=0.80in,bottom=0.82in,columnsep=0.26in,headsep=0.22in,footskip=0.34in]{geometry}
\usepackage{microtype}
\usepackage{graphicx}
\usepackage[table]{xcolor}
\usepackage{booktabs,tabularx,array,multirow}
\usepackage{caption}
\usepackage{enumitem}
\usepackage{titlesec}
\usepackage{fancyhdr}
\usepackage{ragged2e}
\usepackage{stfloats}
\usepackage[numbers,sort&compress]{natbib}
\usepackage[hyphens]{url}
\usepackage[colorlinks=true,linkcolor=accent,citecolor=accent,urlcolor=accent,pdftitle={Performance Before Outcomes},pdfauthor={Jepson Taylor}]{hyperref}

\definecolor{ink}{HTML}{0B0B0B}
\definecolor{ink2}{HTML}{52514E}
\definecolor{accent}{HTML}{1C5CAB}
\definecolor{men}{HTML}{1F5FB0}
\definecolor{women}{HTML}{C2501F}
\definecolor{rulegray}{HTML}{C9C8C2}
\definecolor{band}{HTML}{F4F3EF}
\color{ink}

\titleformat{\section}{\sffamily\bfseries\fontsize{11}{13}\selectfont}{\thesection}{0.6em}{}
\titleformat{\subsection}{\sffamily\bfseries\fontsize{9.4}{11.5}\selectfont}{\thesubsection}{0.5em}{}
\titleformat{\subsubsection}[runin]{\sffamily\bfseries\fontsize{9}{11}\selectfont}{}{0em}{}[.]
\titlespacing*{\section}{0pt}{11pt plus 2pt minus 2pt}{4.5pt}
\titlespacing*{\subsection}{0pt}{8pt plus 2pt minus 1pt}{3pt}
\titlespacing*{\subsubsection}{0pt}{5pt plus 1pt minus 1pt}{0.5em}
\setlength{\parindent}{1.1em}\setlength{\parskip}{0pt plus 0.9pt}

\captionsetup{font={small},labelfont={bf,sf},labelsep=period,justification=justified,singlelinecheck=false,skip=5pt}
\captionsetup[table]{position=top,skip=4pt}
\renewcommand{\topfraction}{0.92}\renewcommand{\dbltopfraction}{0.92}\renewcommand{\bottomfraction}{0.6}
\renewcommand{\textfraction}{0.06}\renewcommand{\floatpagefraction}{0.75}\renewcommand{\dblfloatpagefraction}{0.7}
\setcounter{topnumber}{3}\setcounter{dbltopnumber}{3}\setcounter{totalnumber}{4}
\setlength{\textfloatsep}{10pt plus 2pt minus 2pt}\setlength{\dbltextfloatsep}{10pt plus 2pt minus 2pt}\setlength{\floatsep}{8pt plus 2pt minus 2pt}
\newcolumntype{L}{>{\RaggedRight\arraybackslash}X}
\newcolumntype{P}[1]{>{\RaggedRight\arraybackslash}p{#1}}
\newcommand{\tablenote}[1]{\par\vspace{2pt}{\scriptsize\RaggedRight #1\par}}
\renewcommand{\arraystretch}{1.08}
\setlength{\tabcolsep}{4pt}

\newtheoremstyle{pbo}{5pt plus 1pt}{5pt plus 1pt}{\itshape}{}{\sffamily\bfseries\small}{.}{0.5em}{\thmname{#1}\thmnumber{ #2}\thmnote{ (#3)}}
\theoremstyle{pbo}
\newtheorem{result}{Result}
\newtheoremstyle{pbor}{5pt}{5pt}{}{}{\sffamily\bfseries\small}{.}{0.5em}{}
\theoremstyle{pbor}
\newtheorem*{remark}{Remark}

\pagestyle{fancy}\fancyhf{}
\renewcommand{\headrulewidth}{0.3pt}\renewcommand{\headrule}{\color{rulegray}\hrule height 0.3pt}
\fancyhead[L]{\sffamily\scriptsize\color{ink2}PERFORMANCE BEFORE OUTCOMES}
\fancyhead[R]{\sffamily\scriptsize\color{ink2}Taylor \textperiodcentered{} preprint, November 2026}
\fancyfoot[C]{\sffamily\scriptsize\color{ink2}\thepage}

\newcommand{\E}{\operatorname{E}}
\newcommand{\Var}{\operatorname{Var}}
\newcommand{\Cov}{\operatorname{Cov}}
\newcommand{\ccm}{\,cm$^3$}
\newcommand{\Hs}{\textit{H.~sapiens}}
\newcommand{\Homo}{\textit{Homo}}
\newcommand{\M}[1]{{\color{men}\textbf{#1}}}
\newcommand{\W}[1]{{\color{women}\textbf{#1}}}
\urlstyle{same}
\providecommand{\doi}[1]{doi:\,\href{https://doi.org/#1}{\nolinkurl{#1}}}
\setlist{leftmargin=1.3em,itemsep=1pt,topsep=3pt,parsep=0pt}
\frenchspacing
\emergencystretch=1.2em

% Generated by code/make_numbers.py from results/*.json. Do not edit by hand.
\newcommand{\CcAgeMen}{20.8}
\newcommand{\CcAgeN}{171}
\newcommand{\CcAgeWomen}{14.1}
\newcommand{\CcArrForN}{41}
\newcommand{\CcArrMenN}{148}
\newcommand{\CcArrWomenN}{151}
\newcommand{\CcAutoMen}{31}
\newcommand{\CcAutoWomen}{8}
\newcommand{\CcAutoWomenFor}{12}
\newcommand{\CcBinN}{339}
\newcommand{\CcBodyN}{58}
\newcommand{\CcBodyPlump}{79}
\newcommand{\CcBodySlim}{21}
\newcommand{\CcBrideServicePct}{10}
\newcommand{\CcBrideWealthPct}{57}
\newcommand{\CcDivAdulH}{2}
\newcommand{\CcDivAdulW}{54}
\newcommand{\CcDivCruelH}{46}
\newcommand{\CcDivCruelW}{0}
\newcommand{\CcDivN}{160}
\newcommand{\CcDivSterH}{13}
\newcommand{\CcDivSterW}{29}
\newcommand{\CcDivSuppH}{21}
\newcommand{\CcDivSuppW}{0}
\newcommand{\CcDowryPct}{3.4}
\newcommand{\CcEaN}{1{,}291}
\newcommand{\CcForServicePct}{16.5}
\newcommand{\CcGap}{6.7}
\newcommand{\CcGapCI}{5.9--7.4}
\newcommand{\CcGapHigh}{8.5}
\newcommand{\CcGapLow}{4.6}
\newcommand{\CcGapPolyCI}{0.41--0.75}
\newcommand{\CcGapPolyN}{128}
\newcommand{\CcGapPolyR}{0.60}
\newcommand{\CcGroomPct}{67}
\newcommand{\CcMenOlderPct}{98}
\newcommand{\CcMonoN}{1{,}257}
\newcommand{\CcMonoPct}{16}
\newcommand{\CcPairDau}{44}
\newcommand{\CcPairSame}{96}
\newcommand{\CcPairSon}{2}
\newcommand{\CcParMen}{30}
\newcommand{\CcParWomen}{45}
\newcommand{\CcParWomenFor}{54}
\newcommand{\CcPolyForMed}{9}
\newcommand{\CcPolyForUnder}{74}
\newcommand{\CcPolySccsMed}{5}
\newcommand{\CcProdServicePct}{8.0}
\newcommand{\CcRelMenOlder}{87}
\newcommand{\CcRelN}{70}
\newcommand{\CcRelWomenOlderN}{2}
\newcommand{\CcSccsN}{186}
\newcommand{\CcTrN}{1{,}272}
\newcommand{\CcWivesHunt}{13}
\newcommand{\CcWivesHuntFor}{40}
\newcommand{\CcWivesLeadFor}{45}
\newcommand{\CcWivesRank}{59}
\newcommand{\DbMWBelAtt}{24.9}
\newcommand{\DbMWBelInt}{16.4}
\newcommand{\DbMWDidAtt}{42}
\newcommand{\DbMWDidInt}{11}
\newcommand{\DbMWErrBel}{11.4}
\newcommand{\DbMWErrSaid}{13.4}
\newcommand{\DbMWSaidAtt}{17.9}
\newcommand{\DbMWSaidInt}{20.9}
\newcommand{\DbWMBelAtt}{35.5}
\newcommand{\DbWMBelInt}{12.5}
\newcommand{\DbWMDidAtt}{54}
\newcommand{\DbWMDidInt}{3}
\newcommand{\DbWMErrBel}{10.8}
\newcommand{\DbWMErrSaid}{14.7}
\newcommand{\DbWMSaidAtt}{27.0}
\newcommand{\DbWMSaidInt}{19.4}
\newcommand{\DcMR}{$-$0.22}
\newcommand{\DcWBot}{45}
\newcommand{\DcWR}{$-$0.23}
\newcommand{\DcWTop}{33}
\newcommand{\DrLikeR}{0.12}
\newcommand{\DrMAct}{37}
\newcommand{\DrMExp}{52}
\newcommand{\DrMProj}{0.51}
\newcommand{\DrMR}{0.15}
\newcommand{\DrMatch}{17.1}
\newcommand{\DrMatchInd}{18.3}
\newcommand{\DrPairs}{4{,}011}
\newcommand{\DrWAct}{48}
\newcommand{\DrWExp}{52}
\newcommand{\DrWR}{0.12}
\newcommand{\DriftNe}{1{,}110}
\newcommand{\DriftSD}{2.3}
\newcommand{\DriftZ}{5.9}
\newcommand{\DvMTarget}{20}
\newcommand{\DvWTarget}{18}
\newcommand{\DxMPanel}{5.9}
\newcommand{\DxMSelf}{6.9}
\newcommand{\DxWPanel}{6.5}
\newcommand{\DxWSelf}{7.2}
\newcommand{\FosAus}{455}
\newcommand{\FosAusN}{32}
\newcommand{\FosAusSD}{62}
\newcommand{\FosBreakA}{2.1}
\newcommand{\FosBreakACI}{1.75--2.20}
\newcommand{\FosBreakB}{1.3}
\newcommand{\FosBreakBCI}{1.25--1.80}
\newcommand{\FosDAICLin}{26.4}
\newcommand{\FosDBICLin}{11.5}
\newcommand{\FosDuMatch}{54}
\newcommand{\FosDuN}{94}
\newcommand{\FosDuR}{0.997}
\newcommand{\FosDuSlope}{0.405}
\newcommand{\FosDuSlopeSE}{0.024}
\newcommand{\FosEH}{678}
\newcommand{\FosEHN}{9}
\newcommand{\FosEHSD}{105}
\newcommand{\FosErecEarly}{823}
\newcommand{\FosErecEarlyN}{24}
\newcommand{\FosErecEarlySD}{142}
\newcommand{\FosErecLate}{1{,}063}
\newcommand{\FosErecLateN}{23}
\newcommand{\FosErecLateSD}{205}
\newcommand{\FosFlo}{419}
\newcommand{\FosFloN}{1}
\newcommand{\FosFold}{3.2}
\newcommand{\FosGens}{127{,}600}
\newcommand{\FosHaldHi}{1.16}
\newcommand{\FosHaldLo}{0.88}
\newcommand{\FosHolDrop}{9.3}
\newcommand{\FosMP}{1{,}244}
\newcommand{\FosMPN}{27}
\newcommand{\FosMPSD}{128}
\newcommand{\FosN}{985}
\newcommand{\FosNBins}{26}
\newcommand{\FosNHol}{700}
\newcommand{\FosNRate}{278}
\newcommand{\FosNal}{545}
\newcommand{\FosNalN}{3}
\newcommand{\FosNalSD}{74}
\newcommand{\FosNea}{1{,}386}
\newcommand{\FosNeaN}{30}
\newcommand{\FosNeaSD}{171}
\newcommand{\FosSDbin}{0.132}
\newcommand{\FosSDhol}{0.117}
\newcommand{\FosSapH}{1{,}323}
\newcommand{\FosSapHN}{700}
\newcommand{\FosSapHSD}{155}
\newcommand{\FosSapP}{1{,}458}
\newcommand{\FosSapPN}{136}
\newcommand{\FosSapPSD}{140}
\newcommand{\FosSapSlope}{0.04}
\newcommand{\FosSapSlopeSE}{0.18}
\newcommand{\FosSegA}{0.08}
\newcommand{\FosSegB}{0.74}
\newcommand{\FosSegC}{0.41}
\newcommand{\FosSlopeBin}{0.416}
\newcommand{\FosSlopeBinSE}{0.034}
\newcommand{\FosSlopeSpec}{0.462}
\newcommand{\FosSlopeSpecSE}{0.011}
\newcommand{\FosSpan}{3.2}
\newcommand{\FosThreeBreaks}{2.1, 1.2, 0.7}
\newcommand{\FosTotalSD}{13.3}
\newcommand{\GbAgeGap}{2.7}
\newcommand{\GbMAffair}{21}
\newcommand{\GbMGini}{0.67}
\newcommand{\GbMMean}{15.4}
\newcommand{\GbMMed}{6}
\newcommand{\GbMPaid}{13}
\newcommand{\GbMTopTwenty}{71}
\newcommand{\GbMYearTwo}{19}
\newcommand{\GbMeanRatio}{2.8}
\newcommand{\GbSerialMAbs}{5.7}
\newcommand{\GbSerialMCI}{+2.3 to +9.1}
\newcommand{\GbSerialWAbs}{3.4}
\newcommand{\GbSerialWCI}{$-$6.6 to $-$0.2}
\newcommand{\GbWAffair}{13}
\newcommand{\GbWGini}{0.60}
\newcommand{\GbWMean}{5.5}
\newcommand{\GbWMed}{3}
\newcommand{\GbWPaid}{2}
\newcommand{\GbWTopTwenty}{64}
\newcommand{\GbWYearTwo}{11}
\newcommand{\GeChlMFive}{10.3}
\newcommand{\GeChlMOne}{21.5}
\newcommand{\GeChlMThree}{12.4}
\newcommand{\GeChlWFive}{19.3}
\newcommand{\GeChlWOne}{9.7}
\newcommand{\GeChlWThree}{10.7}
\newcommand{\GeChlWTwo}{8.8}
\newcommand{\GeEduMPair}{+0.014}
\newcommand{\GeFloorM}{7.1}
\newcommand{\GeFloorMCI}{2.5--11.9}
\newcommand{\GeFloorW}{7.6}
\newcommand{\GeFloorWCI}{2.7--12.1}
\newcommand{\GeGateMNew}{31}
\newcommand{\GeGateMOld}{9}
\newcommand{\GeGateWNew}{62}
\newcommand{\GeGateWOld}{18}
\newcommand{\GeIncMConv}{+0.001}
\newcommand{\GeIncMPair}{+0.045}
\newcommand{\GeIncMPairCI}{[+0.037, +0.053]}
\newcommand{\GeIncWConv}{$-$0.037}
\newcommand{\GeIncWConvCI}{[$-$0.047, $-$0.027]}
\newcommand{\GeIncWPair}{0.000}
\newcommand{\GeIncWPairCI}{[$-$0.008, +0.007]}
\newcommand{\GeN}{34{,}086}
\newcommand{\GeNevMFive}{2.6}
\newcommand{\GeNevMOne}{15.0}
\newcommand{\GeNeverMNew}{13.2}
\newcommand{\GeNeverMOld}{4.5}
\newcommand{\GePairIncMNew}{+0.079}
\newcommand{\GePairIncMNewCI}{[+0.056, +0.102]}
\newcommand{\GePairIncMOld}{+0.026}
\newcommand{\GePairIncWNew}{+0.020}
\newcommand{\GePairIncWNewCI}{[+0.002, +0.037]}
\newcommand{\GePairIncWOld}{$-$0.009}
\newcommand{\GePreMPair}{+0.022}
\newcommand{\GeVocMPair}{+0.010}
\newcommand{\GhHusbNew}{33}
\newcommand{\GhHusbOld}{37}
\newcommand{\GhWifeNew}{35}
\newcommand{\GhWifeOld}{32}
\newcommand{\GrMRareChl}{21.1}
\newcommand{\GrMRareTot}{$-$0.126}
\newcommand{\GrMWeeklyChl}{10.5}
\newcommand{\GrMWeeklyTot}{$-$0.054}
\newcommand{\GrWRareChl}{16.3}
\newcommand{\GrWRareTot}{$-$0.208}
\newcommand{\GrWWeeklyChl}{10.8}
\newcommand{\GrWWeeklyTot}{$-$0.130}
\newcommand{\GssChlM}{15.1}
\newcommand{\GssChlW}{12.2}
\newcommand{\GssDelayM}{2.7}
\newcommand{\GssDelayW}{3.0}
\newcommand{\GssEduMNum}{$-$0.081}
\newcommand{\GssEduWNum}{$-$0.126}
\newcommand{\GssIncMTot}{+0.012}
\newcommand{\GssN}{34{,}104}
\newcommand{\GssNFile}{75{,}699}
\newcommand{\GssNM}{15{,}118}
\newcommand{\GssNW}{18{,}986}
\newcommand{\GssNoPNewM}{15.5}
\newcommand{\GssNoPNewW}{11.1}
\newcommand{\GssNoPOldM}{8.7}
\newcommand{\GssNoPOldW}{7.1}
\newcommand{\GssPartManyM}{2.24}
\newcommand{\GssPartManyW}{1.88}
\newcommand{\GssPartOneM}{2.56}
\newcommand{\GssPartOneW}{2.58}
\newcommand{\GssShareNewM}{42}
\newcommand{\GssShareNewW}{36}
\newcommand{\GssShareOldM}{29}
\newcommand{\GssShareOldW}{23}
\newcommand{\GssSpNew}{0.59}
\newcommand{\GssSpOld}{0.62}
\newcommand{\GssViaM}{71}
\newcommand{\GssViaW}{48}
\newcommand{\JurGapMax}{0.04}
\newcommand{\JurKAmb}{31}
\newcommand{\JurKAtt}{6}
\newcommand{\JurKFun}{11}
\newcommand{\JurKInt}{31}
\newcommand{\JurKSin}{34}
\newcommand{\JurOneAmb}{0.11}
\newcommand{\JurOneAtt}{0.41}
\newcommand{\JurOneFun}{0.27}
\newcommand{\JurOneInt}{0.11}
\newcommand{\JurOneSin}{0.10}
\newcommand{\NetS}{2.6}
\newcommand{\SdAge}{26.4}
\newcommand{\SdM}{277}
\newcommand{\SdN}{551}
\newcommand{\SdPopZeroM}{6.5}
\newcommand{\SdPopZeroW}{1.5}
\newcommand{\SdRevAttM}{21.5}
\newcommand{\SdRevAttW}{15.3}
\newcommand{\SdRevDiffAtt}{6.2}
\newcommand{\SdRevDiffCIAtt}{[3.1, 9.1]}
\newcommand{\SdRevDiffCIInt}{[$-$5.9, $-$0.1]}
\newcommand{\SdRevDiffInt}{$-$2.9}
\newcommand{\SdRevIntM}{1.2}
\newcommand{\SdRevIntW}{4.1}
\newcommand{\SdRows}{8{,}378}
\newcommand{\SdSat}{0.15}
\newcommand{\SdSatAtt}{$-$0.06}
\newcommand{\SdSatCI}{0.02--0.29}
\newcommand{\SdSatN}{202}
\newcommand{\SdSelfAtt}{0.29}
\newcommand{\SdSelfFun}{0.28}
\newcommand{\SdSelfInt}{0.05}
\newcommand{\SdSelfSin}{$-$0.01}
\newcommand{\SdSessions}{21}
\newcommand{\SdStAttM}{27.0}
\newcommand{\SdStIntM}{19.4}
\newcommand{\SdStIntW}{20.9}
\newcommand{\SdStSinW}{18.3}
\newcommand{\SdW}{274}
\newcommand{\SdYesM}{47}
\newcommand{\SdYesW}{37}
\newcommand{\ThBreadthTen}{8.6}
\newcommand{\ThDelivErr}{0.004}
\newcommand{\ThFloorTen}{0.19}
\newcommand{\ThFloorThirty}{0.50}
\newcommand{\ThFloorTwenty}{0.35}
\newcommand{\ThIdentErr}{0.0007}
\newcommand{\ThJuryErr}{0.003}
\newcommand{\ThTrack}{0.012}
\newcommand{\ThTrackGen}{81}
\newcommand{\ThTrackRatio}{119}
\newcommand{\ThValOne}{0.30}
\newcommand{\ThValTen}{0.81}

\begin{document}
\onecolumn
\thispagestyle{empty}
\begingroup
\noindent{\sffamily\scriptsize\color{accent}\textbf{RESEARCH ARTICLE \;\textperiodcentered\; PREPRINT, NOT PEER REVIEWED \;\textperiodcentered\; AI-ASSISTED, DATA AND CODE OPEN \;\textperiodcentered\; NOVEMBER 10, 2026}}\par
\vspace{4pt}
\noindent{\sffamily\bfseries\fontsize{26}{29}\selectfont Performance Before Outcomes\par}
\vspace{4pt}
\noindent{\sffamily\itshape\fontsize{14.5}{17.5}\selectfont\color{ink} It takes a village to read a mind, and that is how we got ours\par}
\vspace{3pt}
\noindent{\sffamily\fontsize{9.6}{12}\selectfont\color{ink2} Early assessment, what women and men are chosen for across cultures, and what it
means for the minds we now build\par}
\vspace{4pt}
\noindent{\sffamily\fontsize{10}{12}\selectfont Jepson Taylor\par}
\vspace{3pt}
\noindent{\color{rulegray}\rule{\textwidth}{0.5pt}}\par
\vspace{2pt}
\noindent{\sffamily\bfseries\small Abstract.}\;
Hominin brains tripled in volume in three million years. Selection can act only on what is known when partners, in-laws and
allies decide, yet what matters most, a lifetime of provisioning and children raised to adulthood, is known only decades
later. We propose that repeated public performance (play, contests, shared work, storytelling) solved this timing problem,
and we test what can be tested in four open data sets. In \FosN{} fossil crania the long-run rate of expansion is about
$10^{-4}$ standard deviations per generation, in steps: selection was never the bottleneck, so brain size tracked a moving
optimum. On \SdRows{} four-minute dates two observers agree only \JurOneInt{} about a stranger's intelligence, but agreement
grows with the number of observers exactly as the Spearman--Brown law predicts, so that \JurKInt{} observers give a
dependable verdict; and people act on what one look can read, not on what they say they want. In a world
sample of societies a daughter chooses her spouse alone in \CcAutoWomen\% and a son in \CcAutoMen\%; in \CcAgeN{} forager groups men marry
\CcGap{} years later than women. Among \GssN{} adults in the United States, selection on men acts at pairing, as a floor,
and on women afterwards, as a ceiling. Eight formal results connect the pieces, and the same mathematics governs how
artificial minds are evaluated and selected. No ancestral parameter is estimated; twelve predictions say what would
refute the account.\par
\vspace{3pt}
\noindent{\sffamily\bfseries\small Keywords.}\; {\small human brain evolution; mate choice; social selection; early assessment; general intelligence;
cross-cultural marriage; sex differences in selection; Spearman--Brown; evaluation of AI systems.}\par
\vspace{4pt}
\noindent\includegraphics[width=\textwidth]{figures/fig1_main.pdf}\par
\vspace{-1pt}
\captionof{figure}{\textbf{The argument in five steps, each from open data.} \textbf{(1)}~Endocranial volume of every dated hominin
cranium in the DeSilva et~al.\ compilation~\cite{desilva2021,desilva2023}; line: best two-break trend; grey crosses:
\textit{H.~naledi} and \textit{H.~floresiensis} (Section~\ref{sec:puzzle}). \textbf{(2)}~Agreement between two panels of
observers about the same person after a four-minute meeting, against panel size, in the Columbia speed-dating
experiment~\cite{fisman2006}; ticks mark the panel needed for agreement of 0.8 (Section~\ref{sec:early}).
\textbf{(3)}~How a spouse is chosen for sons and for daughters in the Standard Cross-Cultural
Sample~\cite{broude1983,kirby2016} (Section~\ref{sec:cultures}). \textbf{(4)}~Mean age at first marriage in \CcAgeN{}
hunter-gatherer groups~\cite{binford2001}; bars are means. \textbf{(5)}~Adults aged 45--64 with no children, by fifths of own
income, General Social Survey 1972--2024~\cite{gss2024}; bars are 95\% intervals (Section~\ref{sec:today}). Blue squares
are men and orange circles are women throughout the paper.}
\label{fig:main}
\endgroup
\twocolumn

\section{Introduction}
\label{sec:intro}

Between the australopiths and the people of the last Ice Age, hominin brain volume grew about \FosFold-fold
(Fig.~\ref{fig:main}, panel~1). Brain tissue is expensive: at four to five years of age a child's brain uses glucose at a
rate equal to about two-thirds of the body's resting metabolism~\cite{kuzawa2014}, and the adult bill had to be paid by a
change of diet and gut~\cite{aiello1995,carmody2011}. Most explanations ask why the expense was worth it. Larger brains
helped individuals extract food from difficult environments~\cite{kaplan2000,gonzalezforero2018}, manage many
relationships~\cite{dunbar1998}, learn from others~\cite{muthukrishna2018}, compete with other
people~\cite{flinn2005}, and attract partners and allies~\cite{miller2000book,westeberhard1983,nesse2007}. These accounts
differ about the benefit. They share an unstated assumption: that the benefit was visible to selection.

\subsubsection{The timing problem} For much of what matters in a human life, that assumption fails at the moment
it is needed. Across 40 foraging societies hunting skill peaks at about 33 years~\cite{koster2020}; among Tsimane hunters
caloric return rates peak near 45, some two decades after adult body size is reached~\cite{gurven2006}. How many children a
person raises to adulthood is not known until late middle age. Yet in \CcAgeN{} hunter-gatherer groups women first marry
at \CcAgeWomen{} years on average and men at \CcAgeMen{} (Section~\ref{sec:cultures}). Wherever reproduction depends on
the decisions of other people (whom to marry, whom to give a daughter to, whom to help, whom to follow), selection acting
through those decisions cannot see anything that has not yet been shown. A capacity that pays off at 40 and is invisible
at 20 is, for the purposes of social selection, not there.

\subsubsection{The proposal} Humans fill the gap with performance. Children play at the tasks of adults in front of
other children and adults; adolescents hunt, gather, joke, argue, sing and tell stories where everyone can compare them;
suitors work for their in-laws for months or years before a marriage is settled; evenings by the fire are given to
stories. We argue that this repeated, public, comparable display did real evolutionary work. It let partners, parents and
allies read slow-to-mature capacities years before lifetime outcomes existed, and so allowed their decisions to turn
differences in capacity into differences in reproduction. \emph{Performance before outcomes} is the claim that early
assessment, not only the usefulness of intelligence, shaped how strongly and in which direction human minds were selected.

\subsubsection{What this paper contributes}
\begin{enumerate}
\item \textbf{The puzzle, measured.} The long-run rate of brain expansion was about $10^{-4}$ standard deviations per
generation, in steps separated by stasis (Section~\ref{sec:puzzle}). Even weak selection outruns that a hundredfold, so
the question is what moved the optimum.
\item \textbf{Early reading, measured.} One observer reads a stranger's mind poorly; a jury reads it well, by a law that
fits \SdRows{} encounters to within \JurGapMax{} (Section~\ref{sec:early}). And choice follows legibility: people act on
what an encounter lets them read.
\item \textbf{Eight results} that say exactly how accuracy, juries, fairness, breadth, floors and energy combine
(Section~\ref{sec:theory}), each checked numerically.
\item \textbf{Women and men, across cultures and today.} In the ethnographic record a jury of kin decides, more firmly
for daughters; the groom pays, serves and waits; and marriages end over conduct, not capacity
(Sections~\ref{sec:sexes} and~\ref{sec:cultures}). Among \GssN{} living adults selection on men sits at pairing and
selection on women after it (Section~\ref{sec:today}).
\item \textbf{A bridge to artificial minds.} Evaluating a model before deployment is performance before outcomes, and
the eight results apply to it without change (Box~1).
\item \textbf{Twelve predictions} that could refute the account (Section~\ref{sec:predictions}).
\end{enumerate}

\subsubsection{What it does not claim} Performance is not offered as the only cause of brain expansion; it is a claim about
the \emph{demand} side that complements ecological, social and cultural accounts of why capacity pays. Nothing here implies
that the sexes differ in intelligence: genes spend half their time in each sex, so selection applied through either sex
moves both~\cite{lande1980}. Nothing here ranks living populations or measures anyone's worth; every data set describes
variation among individuals or among customs. And no parameter of any ancestral population is estimated. The fossil record
constrains the rate, living people and recorded societies supply the measurements, and the bridge between them is theory
that Section~\ref{sec:predictions} exposes to test.

\section{The puzzle, measured}
\label{sec:puzzle}

\subsubsection{Data} We use the compilation of endocranial volumes assembled by DeSilva and
colleagues~\cite{desilva2021,desilva2023}: \FosN{} hominin crania from 7 million years ago to the present, each with a
volume and a date. As a check on the file we reproduce the authors' published mean for Pleistocene \Hs{}
(\FosSapP\ccm, $n=\FosSapPN$) exactly, and \FosDuMatch{} specimens shared with the independent compilation of Du
et~al.~\cite{du2018} agree to $r=\FosDuR$ (Methods).

\subsubsection{The rate is slow} Over the last \FosSpan{} million years (\FosNRate{} Pleistocene crania, excluding the two
late small-brained species), log volume rose by \FosSlopeSpec{} $\pm$ \FosSlopeSpecSE{} per million years fitted to
specimens, and by \FosSlopeBin{} $\pm$ \FosSlopeBinSE{} across 100,000-year bins. The binned value coincides with the
\FosSlopeBin{} obtained by Gingerich from an independent series~\cite{gingerich2022}, and the Du et~al.\ data give
\FosDuSlope{} $\pm$ \FosDuSlopeSE. The standard deviation of log volume is \FosSDbin{} within bins and \FosSDhol{} among
\FosNHol{} recent skulls. With 25-year generations the rate is \FosHaldLo--\FosHaldHi$\times10^{-4}$ standard deviations per
generation: about \FosTotalSD{} standard deviations spread over some \FosGens{} generations.

\subsubsection{The pattern is steps and stasis} A constant trend is rejected in favour of a broken one ($\Delta$AICc
$=\FosDAICLin$, $\Delta$BIC $=\FosDBICLin$). The best two-break model has almost no trend before \FosBreakA{} million years
ago (\FosSegA{} per million years; 95\% interval for the break \FosBreakACI), a fast phase to \FosBreakB{} (\FosSegB;
interval \FosBreakBCI) and a slower rise since (\FosSegC). A three-break model (\FosThreeBreaks{} million years ago) adds
a pause between about 1.2 and 0.7 million years ago; the two criteria disagree on whether the data require it.
Table~\ref{tab:fossils} gives the group means. Volume did not rise within \Hs{} during the Pleistocene
(\FosSapSlope{} $\pm$ \FosSapSlopeSE{} per million years), and it is \FosHolDrop\% lower in the last 12,000 years, a
decline whose timing and reality are debated~\cite{desilva2023,villmoare2022}.

\begin{table}[t]
\caption{Endocranial volume by group in the compilation of DeSilva et~al.~\cite{desilva2021,desilva2023} (our tabulation).}
\label{tab:fossils}
\small
\begin{tabularx}{\columnwidth}{@{}Lrrr@{}}
\toprule
Group & $n$ & Mean (cm$^3$) & SD \\
\midrule
Early hominins and australopiths & \FosAusN & \FosAus & \FosAusSD \\
Early \Homo & \FosEHN & \FosEH & \FosEHSD \\
\textit{H.~erectus}, before 1 Ma & \FosErecEarlyN & \FosErecEarly & \FosErecEarlySD \\
\textit{H.~erectus}, after 1 Ma & \FosErecLateN & \FosErecLate & \FosErecLateSD \\
Middle Pleistocene \Homo & \FosMPN & \FosMP & \FosMPSD \\
Neanderthals & \FosNeaN & \FosNea & \FosNeaSD \\
\Hs, Pleistocene & \FosSapPN & \FosSapP & \FosSapPSD \\
\Hs, last 12,000 years & \FosSapHN & \FosSapH & \FosSapHSD \\
\addlinespace[2pt]
\textit{H.~naledi} & \FosNalN & \FosNal & \FosNalSD \\
\textit{H.~floresiensis} & \FosFloN & \FosFlo & -- \\
\bottomrule
\end{tabularx}
\end{table}

\subsubsection{Three consequences} First, drift is not enough. Under neutral drift the mean wanders by
$\sqrt{h^2 t/N_e}$ standard deviations in $t$ generations~\cite{lande1976}. With heritability $h^2=0.4$ and a long-term
effective size of $10^4$~\cite{lidurbin2011} that is \DriftSD{} standard deviations against \FosTotalSD{} observed, a
\DriftZ-fold excess; for a change this large to fall within the 95\% range of drift, the effective size would have to
stay below about \DriftNe{} for three million years.

Second, and more important, selection was never the limiting factor. The observed rate requires a net selection
differential of only \NetS$\times10^{-4}$ standard deviations per generation at $h^2=0.4$: reproduction would have to
favour the larger-brained by one part in four thousand. Section~\ref{sec:theory} shows that the mildest plausible
filter delivers a hundred times that. A trait that can respond a hundred times faster than it did was not being pushed
steadily against resistance. It was sitting near an optimum set by the balance of benefit and cost, and following that
optimum as it moved. This reading agrees with two recent analyses of cranial form: most of the increase happened within
lineages rather than by replacement of species~\cite{puschel2024}, and the shape of the braincase across eight groups of
\Homo{} fits stasis and an unbiased random walk better than steady directional change~\cite{hubbe2026}. The thing to
explain is what moved the optimum upward in steps, and kept it up.

Third, the optimum did not rise everywhere. \textit{Homo naledi} (\FosNal\ccm) and \textit{H.~floresiensis}
(\FosFlo\ccm) lived at the same time as large-brained people, yet their own brains were the size of
australopiths'~\cite{berger2015,dirks2017,sutikna2016}. Whatever raised the optimum was a feature of particular ways of
life, not of the genus.

\subsubsection{A blunt proxy} Endocranial volume is what the fossils give us; it is not the trait of interest. Among
living people brain volume is only weakly related to test scores: a meta-analysis of 148 samples found a robust but weak
association~\cite{pietschnig2015}, and the correlation is 0.19 among 13,608 adults~\cite{nave2019}. The internal
organization of the brain changed on its own schedule~\cite{poncedeleon2021,neubauer2018}. We use volume because it is the
only long-run measure there is, and make no claim that it measures the capacities that matter well.

\section{Reading a life early}
\label{sec:early}

\begin{figure*}[t]
\centering
\includegraphics[width=\textwidth]{figures/fig2_early.pdf}
\caption{\textbf{Evidence before outcomes.} \textbf{(A)}~Published values on one lifeline. Violet: hunting skill as a share of
its peak in a synthesis of 23,747 records from 1,821 foragers at 40 sites~\cite{koster2020}. Green: ages at which
upper-body strength, kill rate and caloric return rate peak among Tsimane hunters~\cite{gurven2006}. Red: scores on the same
reasoning test at 11 and 77 years correlate 0.63~\cite{deary2000}. Pink: the age (4--5 years) at which the brain's glucose
use is largest relative to resting metabolism~\cite{kuzawa2014}. Orange and blue bands: mean age at first marriage of women
and men in \CcAgeN{} forager groups~\cite{binford2001} (our tabulation). \textbf{(B)}~For five traits, the agreement of two
observers after one four-minute meeting (horizontal) against the trait's weight in the decision to meet again (vertical;
Table~\ref{tab:sd}), speed-dating data~\cite{fisman2006}.}
\label{fig:early}
\end{figure*}

\subsection{Outcomes arrive late; evidence arrives early}

Figure~\ref{fig:early}A places published measurements on a single lifeline. The outcomes that selection ultimately counts
(completed family size, the number of children who survive, a lifetime's production) exist only on the right-hand side of
the figure. Decisions about partners are made on the left. What is available on the left is evidence, and it is
better than one might expect. The average hunter has 89\% of his eventual peak skill at 18, and skill then stays above
that level until 56~\cite{koster2020}; differences among adults are large (some have half the skill of others in the same
community). A reasoning test taken at 11 correlates 0.63 with the same test taken at 77, and 0.73 after correction for the
narrowed range of the survivors~\cite{deary2000}. Table~\ref{tab:early} collects these and other early readings.

\begin{table*}[t]
\caption{What can be read early, and how well. Rows marked \emph{this paper} are computed from the speed-dating
data~\cite{fisman2006}; ``agreement'' is the correlation between two independent observers of the same person.}
\label{tab:early}
\small
\begin{tabularx}{\textwidth}{@{}P{0.17\textwidth}P{0.24\textwidth}Lp{0.14\textwidth}@{}}
\toprule
Capacity & Early evidence & How good is the reading & Source \\
\midrule
Hunting and foraging skill & Returns in adolescence, seen by the whole camp & 89\% of peak skill at 18; peak at 33; adults in one community differ up to twofold & \cite{koster2020} \\
General reasoning & One group test at age 11 & $r=0.63$ with the same test at 77 (0.73 corrected for range) & \cite{deary2000} \\
Intelligence, from behaviour & Strangers watch 15 short videotaped tasks & $r\approx0.5$ with test scores once ratings are pooled over tasks and judges & \cite{borkenau2004} \\
Intelligence, from one meeting & Four minutes, one observer & Agreement \JurOneInt; \JurKInt{} observers needed for 0.8 & this paper \\
Looks & Four minutes & Agreement \JurOneAtt; \JurKAtt{} observers for 0.8 & this paper \\
Humour (``fun'') & Four minutes & Agreement \JurOneFun; \JurKFun{} observers for 0.8 & this paper \\
Ambition; sincerity & Four minutes & Agreement \JurOneAmb{} and \JurOneSin; \JurKAmb{} and \JurKSin{} observers for 0.8 & this paper \\
Physical strength & Photographs of body or face & Judged accurately by strangers & \cite{sell2009} \\
Character and drive & Ratings by acquaintances & Predict school and job performance better than self-ratings & \cite{connelly2010} \\
Storytelling & Reputation among camp-mates & Named storytellers are chosen as camp-mates about twice as often (odds ratio 1.95) and have 0.53 more living children & \cite{smith2017} \\
\bottomrule
\end{tabularx}
\end{table*}

\subsection{The jury law, measured}
\label{sec:jury}

How much can be read in a single short encounter, and how quickly do observers converge? The Columbia speed-dating
experiment~\cite{fisman2006} is close to an ideal instrument: \SdN{} graduate students (\SdW{} women, \SdM{} men, mean age
\SdAge) met in \SdSessions{} sessions; each pair talked for four minutes; each person then rated the other on six
attributes and said yes or no to meeting again. A person was typically rated by about 16 others. For each attribute we
drew two disjoint sets of $k$ raters of the same person, averaged each set, and correlated the two averages across
persons (Methods).

Two lone observers agree \JurOneAtt{} about looks, \JurOneFun{} about how much fun someone is, and only \JurOneInt,
\JurOneAmb{} and \JurOneSin{} about intelligence, ambition and sincerity. (These participants were selected for
academic ability, which narrows the range and lowers agreement about intelligence; the ordering of the traits is the
robust finding.) One meeting is close to a guess about a stranger's mind. But agreement climbs with $k$ exactly as the
Spearman--Brown law requires~\cite{spearman1910,brown1910}:
\begin{equation}
r_k=\frac{k\,r_1}{1+(k-1)\,r_1},
\label{eq:sb}
\end{equation}
with a largest gap of \JurGapMax{} between measured and predicted values across five attributes and $k=1,\dots,8$
(Fig.~\ref{fig:main}, panel~2). Solving for $r_k=0.8$ gives $k=4(1-r_1)/r_1$: \JurKAtt{} observers for looks, \JurKFun{} for
fun, and \JurKInt, \JurKAmb{} and \JurKSin{} for intelligence, ambition and sincerity. Foraging bands hold a few dozen
people~\cite{hill2011}. A lone suitor cannot read a mind in an evening. A camp that has watched someone for years can.

Two further facts sharpen the point. People cannot simply be asked: self-rated intelligence correlates \SdSelfInt{} with
the panel's verdict and self-rated sincerity \SdSelfSin{} (looks \SdSelfAtt, fun \SdSelfFun). And the panel's verdict is
not mere consensus: its rating of intelligence correlates \SdSat{} (95\% interval \SdSatCI; $n=\SdSatN$) with the median
entrance-test score of the college each person had attended, while its rating of looks does not ($r=\SdSatAtt$). In a
sample as selected as an elite graduate school that is a weak signal, but it has the right sign and is specific to the
right trait.

\subsection{People act on what they can read}
\label{sec:legible}

Before the event the same participants divided 100 points among six attributes to say what they wanted. Women gave most
to intelligence (\SdStIntW) and sincerity (\SdStSinW); men to attractiveness (\SdStAttM) and then intelligence
(\SdStIntM). What they did was different (Table~\ref{tab:sd}). For both sexes a yes followed attractiveness above all,
then shared interests and fun; intelligence earned a little, sincerity and ambition nothing.
Figure~\ref{fig:early}B shows why. The weight a trait carried in the decision rises with how well a single meeting reads
it: looks (agreement \JurOneAtt), then fun (\JurOneFun), then intelligence, ambition and sincerity (about 0.1). The
participants said that mind and character mattered most and then could not act on the preference, because four minutes
do not reveal them. A preference is only as strong as the instrument that serves it.

These data cannot fully separate legibility from taste; looks may simply matter more for a first date, and in other
samples measured intelligence was detected by observers without being found attractive~\cite{driebe2021}. Prediction P8
is the test. But the implication for selection does not depend on that. A society in which mates are chosen in an
evening selects on what an evening shows. A society that selects on mind and character must own instruments that read
them: long acquaintance, many observers, and staged, repeated, comparable performance.

\subsection{Why performance beats waiting}

A lifetime outcome is one late observation. A performance has five properties that a lifetime outcome lacks. It is
\emph{early}: it can be staged long before the outcome exists. It is \emph{repeatable}, so that $k$ in
Eq.~\eqref{eq:sb} can be made large. It is \emph{public}, so that many observers see the same event and errors average
out. It is \emph{fair}: everyone attempts the same task under the same conditions among people of the same age, which
strips out differences of luck and inheritance (Result~\ref{res:fair}). And it is \emph{safe}: a contest that does not
kill the loser can be held again tomorrow.

Human societies are full of institutions with these properties. Children in foraging societies learn subsistence skills
largely through play and through practice in mixed-age groups, in public~\cite{lewlevy2017}; athletic contests are
found across the ethnographic record~\cite{deaner2013}. Bride service, the prevailing custom in \CcForServicePct\% of
foraging societies (Section~\ref{sec:cultures}), is an audition lasting months or years in front of the people who
decide. Among the Ju/'hoansi, 81\% of firelit evening conversation was storytelling against 6\% by day, when talk turned
on economic matters and complaints~\cite{wiessner2014}; among the Agta a reputation for storytelling predicted being
chosen as a camp-mate more strongly than a reputation for hunting or fishing, and skilled storytellers had more
surviving children~\cite{smith2017}. Each of these practices turns a private, slow-to-mature capacity into a public,
early, repeated reading.

\section{Eight results}
\label{sec:theory}

The argument can be stated exactly. Let $a$ be an individual's heritable capacity (breeding value, variance
$\sigma_A^2$), $z=a+e$ the expressed capacity, and $h^2=\sigma_A^2/\sigma_z^2$. By the time decisions are made, observers
hold a score $Y$ built from whatever they have seen. Relative fitness is $w$ (mean 1) and $I=\Var(w)$ is the
opportunity for selection~\cite{crow1958}. The change in mean breeding value in one generation is
$\Delta\bar a=\Cov(a,w)$~\cite{price1970}. Proofs are in Appendix~A; every statement was also checked by simulation
(Methods).

\begin{result}[Accuracy $\times$ alignment $\times$ opportunity]
\label{res:identity}
Suppose decisions use only what observers have seen ($w$ is independent of $a$ given $Y$) and the regression of $a$ on
$Y$ is linear. Then, for any rule that turns scores into reproduction,
\begin{equation}
\Delta\bar a=\sigma_A\;\rho_{aY}\;\rho_{Yw}\;\sqrt{I}.
\label{eq:identity}
\end{equation}
If the score reflects breeding value only through expressed capacity, with reliability $R$, then $\rho_{aY}=h\sqrt{R}$.
\end{result}

\noindent The three factors are separate things. \emph{Accuracy} $\rho_{aY}$ is how well the early reading tracks
heritable capacity. \emph{Alignment} $\rho_{Yw}$ is how faithfully the people who decide act on the reading.
\emph{Opportunity} $\sqrt{I}$ is how unequal reproduction is. Information changes nothing where reproduction is equal, and
inequality selects for nothing where it is unrelated to capacity. Two corollaries carry the thesis of the paper.
\emph{Timing}: only evidence available when the decision is made enters $Y$; an outcome that arrives later has
$\rho_{aY}=0$ for that decision however informative it eventually becomes. \emph{Substitution}: any practice that raises
$\rho_{aY}$ at the age of decision is equivalent, in its effect on selection, to a proportional rise in reproductive
inequality, without the conflict that inequality brings.

\begin{result}[Juries and repetition buy accuracy]
\label{res:jury}
If single readings have reliability $r_1$ and independent errors, the mean of $k$ readings has reliability
$r_k=kr_1/(1+(k-1)r_1)$, and $4(1-r_1)/r_1$ readings reach 0.8. If errors are correlated $c$ across readings the
ceiling is $r_1/(r_1+c\,(1-r_1))$.
\end{result}

\noindent The first part is the law that the speed-dating data obey. It does not matter whether the $k$ readings are $k$
observers of one event or one observer of $k$ events: publicity and repetition are the same resource. The second part is
a warning. Readings passed on as gossip share their errors, and shared error cannot be averaged away: with $r_1=0.11$
and $c=0.3$ no number of observers exceeds 0.29. Reputation is cheap but capped. Fresh performances in front of the
people who will decide are what lift the cap.

\begin{result}[Fair contests buy heritable accuracy]
\label{res:fair}
Let a cue be $Y=a+u+\varepsilon$, where $u$ is circumstance (inherited wealth, help from kin, luck, age) and
$\varepsilon$ is noise of observation. Then $\rho_{aY}^2=\sigma_A^2/(\sigma_A^2+\sigma_u^2+\sigma_\varepsilon^2)$, and a
noisy cue beats a precise one whenever it removes more variance of circumstance than it adds noise.
\end{result}

\noindent A lifetime outcome is measured almost without error ($\sigma_\varepsilon^2\approx0$) but is soaked in
circumstance. A footrace, a riddle or an archery match among age-mates is a small controlled experiment: same task, same
conditions, same day. It measures less, but more of what it measures is the person. This is why a brief, rule-bound
display can be a better guide to heritable capacity than the outcome it anticipates, and why the rules, and their
enforcement by onlookers, matter.

\begin{result}[Display shows capacity; only conduct shows delivery]
\label{res:delivery}
Let a partner be worth $v=p\,b$ to a chooser: capacity $b$ times the share $p$ that is delivered, with $b$ and $p$
independent. (i)~$\Cov(b,v)=\bar p\,V_B$ and $\Cov(p,v)=\bar b\,V_P$: capacity is selected in proportion to average
delivery, and delivery in proportion to average capacity. (ii)~If a share $d$ of $n$ observed occasions can be recognized
as tests, capacity is read with reliability $V_B/(V_B+\sigma^2/(dn))$ and delivery with reliability
$V_P/(V_P+\nu/((1-d)n))$. (iii)~After $m$ unrecognized occasions without a default, the default rate is bounded only by
about $3/m$.
\end{result}

\noindent A test shows what a person can do; it cannot show what the person will do once the test is over. Capacity is
therefore read from staged occasions, where people do their best, and reliability from unstaged ones, where they do not
know they are watched; maximum and typical performance correlate only moderately even at work~\cite{sackett1988,beus2012}.
Three things follow. Capacity is worth choosing only where it is delivered, so in a species in which mates and allies
share food and care, capacity and reliability are selected together. Reliability can never be read in an evening:
it takes long acquaintance, life in a group small enough that conduct is seen daily, or an audition such as bride service.
And it is never established, only not yet refuted, which is why assessment continues after the wedding
(Section~\ref{sec:cultures}).

\begin{result}[Breadth aims selection at general ability]
\label{res:breadth}
Let performance in each of $K$ domains be $x_k=g+s_k+\varepsilon_k$, with a general factor $g$ (variance $V_G$) and
independent specific skills $s_k$ (variance $V_S$), and let observers score the average. Then the selection differential on
$g$ relative to that on any one specific skill, each in its own standard deviations, is $K\sqrt{V_G/V_S}$, and the
validity of the score for $g$ is $V_G/(V_G+(V_S+V_\varepsilon)/K)$.
\end{result}

\noindent Someone judged as a hunter, tracker, toolmaker, storyteller, negotiator and parent is being measured, in
effect, on what those tasks have in common. With ten domains, selection aims \ThBreadthTen{} times more sharply at the
general factor than at any specialty in our simulation (Fig.~\ref{fig:model}B), and the validity of the composite rises
from \ThValOne{} to \ThValTen. Breadth of performance is thus a second route to accuracy, and it explains the direction
of human cognitive evolution as well as its amount: toward a general-purpose mind, which is also what comparative work
finds varies among primates~\cite{reader2011,burkart2017}. The converse is a prediction. Where a division of labour lets
each person be judged in one specialty, the premium on general ability falls.

\begin{figure*}[t]
\centering
\includegraphics[width=\textwidth]{figures/fig3_model.pdf}
\caption{\textbf{The model (synthetic; no empirical data).} \textbf{(A)}~Response per generation when the lowest-scoring
fraction of one sex leaves no offspring, for three accuracies of the early reading (Result~\ref{res:floor}); the dashed
line is the long-run rate in the fossil record (Section~\ref{sec:puzzle}). \textbf{(B)}~Selection on the general factor
relative to any one specific skill as the number of judged domains grows: simulation (points) and
Result~\ref{res:breadth} (line), with $V_G=0.3$, $V_S=0.4$. \textbf{(C)}~The favoured brain size under
Result~\ref{res:optimum} as the early reading becomes more accurate, at two costs of tissue.}
\label{fig:model}
\end{figure*}

\begin{result}[A floor is enough]
\label{res:floor}
If the lowest-scoring fraction $q$ leaves no offspring and the rest reproduce equally, the selection intensity is
$i(q)=\phi(z_q)/(1-q)$ standard deviations of the score: \ThFloorTen, \ThFloorTwenty{} and \ThFloorThirty{} for $q=0.1$,
$0.2$ and $0.3$. Applied to one sex, the response is $\tfrac12 h^2\sqrt{R}\,i(q)$ phenotypic standard deviations per
generation.
\end{result}

\noindent Selection does not need a podium. With $h^2=0.4$, a poor reading ($R=0.1$) and one man in ten left without
children, the response is \ThTrack{} standard deviations per generation: \ThTrackRatio{} times the long-run rate in the
fossil record, or one standard deviation in \ThTrackGen{} generations (Fig.~\ref{fig:model}A). Even a glimpse ($R=0.01$)
outruns the record. Exclusion at the bottom is also the socially cheapest form of selection, since it needs no despot
and no harem, only that a few are judged unready and passed over.

\begin{result}[Reading and energy multiply]
\label{res:optimum}
Let capacity grow with the logarithm of brain size $B$, let fitness rise with capacity at rate
$\beta=\beta_{\mathrm{own}}+\rho\,\beta_{\mathrm{social}}$, where $\rho$ is the accuracy of the early reading, and let
tissue cost $\kappa$ per unit: $\ln w=a\beta\ln B-\kappa B$. Then the favoured size is
\begin{equation}
B^{*}=\frac{a\,(\beta_{\mathrm{own}}+\rho\,\beta_{\mathrm{social}})}{\kappa}.
\label{eq:optimum}
\end{equation}
\end{result}

\noindent Because populations can close a gap to the optimum within a few hundred generations
(Result~\ref{res:floor}), the fossil curve is in effect a record of $B^*$ through time. Equation~\eqref{eq:optimum} says
what can move it. Cheaper energy lowers $\kappa$; cooking is the standard candidate~\cite{carmody2011,zohar2022}. A harder or
richer niche raises $\beta_{\mathrm{own}}$~\cite{kaplan2000,gonzalezforero2018}. And anything that makes capacity legible
at the age of decision raises $\rho$, converting social rewards for capacity, which may always have been on offer, into
selection. The terms multiply: better reading is worth most where energy is cheap, and cheap energy is worth most where
capacity can be read (Fig.~\ref{fig:model}C). The functional forms are illustrative, but the comparative statics hold for
any concave benefit and convex cost. Stasis is an unchanged ratio; a step is a change in one of the terms; and a
small-brained species beside a large-brained one is a population in which the ratio fell.

\begin{result}[Pairing, conversion and number]
\label{res:episodes}
Let $z$ be a standardized trait and $c$ relative family size, and let $p$ be the fraction who have any children. Then
\begin{equation}
S=\underbrace{\E[z\mid c>0]}_{\text{access}}+\underbrace{\frac{\Cov(z,c\mid c>0)}{\E[c\mid c>0]}}_{\text{number}},
\qquad
I=\frac{1-p}{p}+\frac{I_{\mathrm{parents}}}{p},
\label{eq:episodes}
\end{equation}
and access splits again into \emph{pairing}, $\E[z\mid\text{ever paired}]$, and \emph{conversion},
$\E[z\mid c>0]-\E[z\mid\text{ever paired}]$.
\end{result}

\noindent The selection differential splits exactly into who is accepted as a partner, who then becomes a parent, and
how many children parents have, and the opportunity for selection splits the same way~\cite{arnold1984}. The split
matters because the episodes are decided by different people on different evidence. Pairing is settled early, by
partners and in-laws reading performance. Conversion and number are settled later, by health, help, luck and choice.
The next three sections show that the first episode has weighed more on men and the later ones on women, in the
ethnographic record and in living people.

\section{What women and men are chosen for}
\label{sec:sexes}

Mate selection is not one event. It is at least five questions: who judges, what they watch, how well they read it, at
which episode of reproduction the judgment bites, and what the judgment rewards. The answers differ by sex
(Table~\ref{tab:sexes}).

\begin{table*}[t]
\caption{The anatomy of selection by sex. Upper rows: the forager and small-scale record, from the cited sources and from
our tabulations of coded ethnographies (Section~\ref{sec:cultures}); these are tendencies with many exceptions. Lower
rows: the United States today (Section~\ref{sec:today}).}
\label{tab:sexes}
\small
\begin{tabularx}{\textwidth}{@{}P{0.13\textwidth}LL@{}}
\toprule
 & \W{Women} & \M{Men} \\
\midrule
Episode that bites & \emph{Conversion and number.} Most women marry and bear children; differences arise mainly in how many
are born and how many survive & \emph{Pairing.} Men's reproductive success usually varies more than
women's~\cite{brown2009,betzig2012}; status raises it mainly through marriage and fertility~\cite{vonrueden2016} \\
Who judges & Her parents, who choose her husband in \CcParWomen\% of societies and must approve in most others; then her
husband and his kin; and the helpers she can recruit~\cite{sear2008,hrdy2009} & The bride's parents and kin; the bride; and
the men whose support he needs~\cite{boehm1999}. He chooses alone in \CcAutoMen\% of societies \\
What they watch & Hard work and foraging returns, fertility and health, character~\cite{marlowe2004}; skill at keeping
children alive and at holding a network of help together & Hunting and provisioning~\cite{smith2004,marlowe2004};
intelligence~\cite{marlowe2004}; generosity and reliability; strength under control~\cite{lidborg2022}; standing among
other men~\cite{vonrueden2016}; storytelling and knowledge~\cite{smith2017} \\
When & First marriage at \CcAgeWomen{} years in foragers; daily work is public from childhood & First marriage at
\CcAgeMen; skill is at 89\% of its peak by 18~\cite{koster2020}; the groom's side pays or serves in \CcGroomPct\% of
societies \\
What ends the marriage & Adultery, leaving for a lover, sterility, laziness, old age (Fig.~\ref{fig:cultures}D) & Cruelty,
taking another wife, failure to provide, desertion, neglect (Fig.~\ref{fig:cultures}D) \\
Size of the effect & Loss of the mother sharply lowers a young child's survival; maternal grandmothers tend to raise it;
fathers do so in only a third of 45 studies~\cite{sear2008} & Status and reproductive success correlate 0.19 across 288
results from 33 societies, 0.18 among foragers~\cite{vonrueden2016}; strength 0.14, rising to 0.16 in high-fertility
samples~\cite{lidborg2022} \\
\midrule
Today: childless, 45--79 & \GssChlW\% & \GssChlM\% \\
Today: pairing, income & \GeIncWPair{} SD (no selection at the door) & \GeIncMPair{} SD (higher earners more often
marry) \\
Today: conversion, income & \GeIncWConv{} SD (higher earners less often become mothers) & \GeIncMConv{} SD \\
Today: number, schooling & \GssEduWNum{} SD & \GssEduMNum{} SD \\
\bottomrule
\end{tabularx}
\end{table*}

\subsection{Men: chosen at the door}

For men the decisive episode has been access. In a meta-analysis of 288 results from 33 non-industrial societies, men's
status (measured as hunting skill, wealth, political influence or physical formidability) correlated 0.19 with
reproductive success, with no difference between foragers and farmers and no difference among the kinds of status; the
effect ran through fertility and number of mates more than through child survival~\cite{vonrueden2016}. Good hunters had
more children in each of the foraging societies reviewed by Smith~\cite{smith2004}, and Hadza women say they place great
weight on a man's foraging and value intelligence more than Hadza men do~\cite{marlowe2004}. In a Finnish population with
every episode on record, variance in mating success explained most of men's excess variance in
fitness~\cite{courtiol2012}. Genetic evidence agrees in kind: mitochondrial lineages imply an effective population of
women more than twice that of men over human history, and up to 17 times during a bottleneck in male lineages 8,000 to
4,000 years ago~\cite{karmin2015}; and in the UK Biobank, carriers of rare damaging variants are more often childless,
men especially, apparently for want of a partner~\cite{gardner2022}.

Three features of this selection are easy to miss. First, the judges are a jury. Where parents and kin arrange or must
approve marriages, as in most societies (Section~\ref{sec:cultures}), a young man is assessed by older people who have
watched him for years and have no romantic stake: exactly the many-observer, many-occasion reading that
Results~\ref{res:jury} and~\ref{res:delivery} say is needed for intelligence and reliability, which one meeting cannot
read. Parents and offspring indeed want different things: across cultures parents put more weight on a prospective
in-law's background, investment and cooperativeness, and offspring more on attractiveness~\cite{buunk2008}. Second, the
filter is a floor, not a podium. In forager groups a median of \CcPolyForMed\% of men have more than one wife, and under
a fifth do in \CcPolyForUnder\% of groups; reproduction is not monopolized by a few. Selection comes from the men who are
judged not ready, not able or not safe, and passed over or made to wait, and Result~\ref{res:floor} shows that this is
ample. Third, force is rewarded only within limits. In a meta-analysis of 96 studies, bodily strength was the one
masculine trait that predicted reproduction ($r=0.14$), while voice pitch, height, testosterone and facial masculinity
did not~\cite{lidborg2022}. The evidence that killers out-reproduce is mixed: positive among the Yanomam\"o and
Nyangatom~\cite{chagnon1988,glowacki2015}, negative among the Waorani~\cite{beckerman2009}. Foragers are known for
coalitions that cut down bullies and boasters~\cite{boehm1999,wrangham2019}, and deference in such societies goes to the
skilled and generous as well as the formidable~\cite{henrich2001}. Cruelty is the commonest failing for which husbands
are divorced (Fig.~\ref{fig:cultures}D). What the jury rewards is closer to competence under control than to dominance.

\subsection{Women: chosen for what follows}

For women the door is rarely closed; the decisive episodes come after. A woman's reproductive success in a
natural-fertility population depends on how many of her children live, and that depends heavily on other people. Across
45 studies the loss of the mother sharply reduced a young child's chance of survival, maternal grandmothers tended to
improve it, and fathers made a detectable difference in only a third of the populations~\cite{sear2008}. Human mothers
raise overlapping, slow-growing, expensive children that no ape mother could support alone, and they do it with
recruited help~\cite{hrdy2009,hawkes1998,kramer2010}. Among the Tsimane a mother's knowledge of plants predicts her
children's health~\cite{mcdade2007}, and her political influence predicts their growth and freedom from
illness~\cite{alami2020}.

That places a specific and heavy demand on women's minds: reading who can be trusted with a child, keeping allies,
managing obligations across households, planning births and food across years. It also means that women are assessed
continuously by the people whose help they need, on evidence that is public and early: who works, who shares, who
keeps the peace, whose children thrive. Men's stated preferences fit. Hadza men place considerable weight on a wife being
hard-working, value fertility more than women do, and, like women, name good character as important~\cite{marlowe2004}.
None of this is a lesser form of selection than the men's. In terms of Eq.~\eqref{eq:episodes} it acts on the later terms
rather than the first, and much of it needs no one's choice at all. Because the same genes build both sexes' brains,
the two pressures add.

\subsection{What the sexes share}

The differences are of emphasis. Across 45 countries both sexes rate kindness, intelligence and health highly in a
long-term partner~\cite{walter2020}; the well-known differences, first documented across 37 cultures~\cite{buss1989}
(men weigh youth and looks more, women weigh resources and standing more), sit on top of that shared core. When people
have a limited budget of mate value to spend they buy sufficiency in the essentials before any luxury~\cite{li2002},
which is floor-setting again. And because choice is mutual in a pair-bonding species, each sex is both the performer and
the jury.

\section{Across cultures: who chooses, who pays, who waits}
\label{sec:cultures}

\begin{figure*}[t]
\centering
\includegraphics[width=\textwidth]{figures/fig4_cultures.pdf}
\caption{\textbf{The cross-cultural record} (our tabulations of coded ethnographies in D-PLACE~\cite{kirby2016}).
\textbf{(A)}~Prevailing transaction at marriage in the Ethnographic Atlas~\cite{murdock1967}; blue: the groom's side pays
or serves; orange: the bride's side pays. \textbf{(B)}~Share of societies in which additional wives go with rank or
wealth, leadership, hunting skill or shamanship, for all societies and for foragers~\cite{white1986}.
\textbf{(C)}~Years by which men marry later than women against the share of men with more than one wife, in
hunter-gatherer groups~\cite{binford2001}. \textbf{(D)}~Number of societies in which each ground for divorce is recorded,
by the spouse at fault~\cite{betzig1989}; cases coded ``both'' or ``unspecified'' are omitted.}
\label{fig:cultures}
\end{figure*}

If the argument is right, the world's marriage customs should look like machinery for reading people early and for
holding them to what was read. Three coded samples of the ethnographic record let us check: the Ethnographic Atlas
(\CcEaN{} societies~\cite{murdock1967}), the Standard Cross-Cultural Sample (\CcSccsN{} societies chosen to be
relatively independent~\cite{murdock1969}) and Binford's \CcBinN{} hunter-gatherer groups~\cite{binford2001}, all in the
open D-PLACE database~\cite{kirby2016}. Societies share ancestry and neighbours, so we report counts and shares and
make no significance claims across them; the codes are ethnographers' judgments, mostly of the nineteenth and
twentieth centuries. Five regularities stand out, and one caution.

\subsubsection{A jury decides, and more firmly for daughters} In the codes of Broude and Greene~\cite{broude1983}, a
son chooses his spouse alone in \CcAutoMen\% of \CcArrMenN{} societies and a daughter in \CcAutoWomen\% of \CcArrWomenN.
Parents choose for \CcParMen\% of sons and \CcParWomen\% of daughters; in the rest the individual chooses but elders must
approve, or both customs exist (Fig.~\ref{fig:main}, panel~3). Where both sexes are coded, daughters are less free than
sons in \CcPairDau{} societies, as free in \CcPairSame, and freer in \CcPairSon. Foragers are no exception: among
\CcArrForN{} foraging societies daughters choose alone in \CcAutoWomenFor\% and parents choose in \CcParWomenFor\%, which
agrees with an independent coding of 190 foraging societies in which marriage was arranged by parents or kin in
85\%~\cite{apostolou2007,walker2011}. Almost everywhere, then, the verdict on a suitor is given by people who have had
years to watch.

\subsubsection{The groom pays or serves} Of \CcTrN{} societies with data, the groom or his kin transfer wealth to the
bride's family in \CcBrideWealthPct\% and the groom works for them in \CcBrideServicePct\%; dowry, the reverse flow,
prevails in \CcDowryPct\% (Fig.~\ref{fig:cultures}A). Among foragers, who have little wealth to transfer, bride service
is twice as common as among food producers (\CcForServicePct\% against \CcProdServicePct\%). It is the man who must
qualify, and where he cannot pay he performs.

\subsubsection{Men wait} In \CcAgeN{} hunter-gatherer groups women first marry at a mean of \CcAgeWomen{} years and men at
\CcAgeMen: \CcGap{} years later (95\% interval \CcGapCI), and later than women in \CcMenOlderPct\% of groups
(Fig.~\ref{fig:main}, panel~4). In the cross-cultural sample men are the older partner at first marriage in
\CcRelMenOlder\% of \CcRelN{} societies and women in \CcRelWomenOlderN{} of them~\cite{whyte1978}. In the United States in
the twentieth century the gap was \GbAgeGap{} years (Section~\ref{sec:today}). The wait is the audition. It falls in
exactly the years when a hunter's skill is still climbing toward its peak (Fig.~\ref{fig:early}A), so a man is judged on
several more years of public record than a woman is. A woman married in her mid-teens cannot have been chosen on a
long record of her own, which is one reason her parents choose and why her assessment continues afterward.

\subsubsection{Polygyny is rationed by the wait, and goes to rank} Only \CcMonoPct\% of \CcMonoN{} societies are
monogamous by rule, but plural marriage is a minority practice within them: the median share of men with more than one
wife is \CcPolyForMed\% across forager groups and \CcPolySccsMed\% across the cross-cultural sample. Where polygyny is
commoner, men wait longer: across \CcGapPolyN{} forager groups the age gap correlates \CcGapPolyR{} (interval
\CcGapPolyCI) with the share of polygynous men, \CcGapLow{} years below the median and \CcGapHigh{} above
(Fig.~\ref{fig:cultures}C). The floor of Result~\ref{res:floor} is thus largely a queue: most men marry in the end, but
some much later, and how long each waits is decided by those watching. Who gets additional wives changes with the
economy. Among foragers extra wives go to skilled hunters in \CcWivesHuntFor\% of societies and to leaders in
\CcWivesLeadFor\%; across all societies they go with rank or wealth in \CcWivesRank\% and with hunting skill in
\CcWivesHunt\% (Fig.~\ref{fig:cultures}B). Once wealth can be stored and inherited, the door rewards circumstance as
well as capacity, which by Result~\ref{res:fair} makes it a worse instrument.

\subsubsection{The audition never ends} Betzig's codes record the grounds on which marriages are dissolved in
\CcDivN{} societies~\cite{betzig1989}. Husbands are divorced for cruelty (\CcDivCruelH{} societies, against
\CcDivCruelW{} for wives), for taking another wife, for failing to provide (\CcDivSuppH{} against \CcDivSuppW), for
desertion and for neglect. Wives are divorced for adultery (\CcDivAdulW{} societies, against \CcDivAdulH{} in which only
the husband's adultery is grounds), for sterility (\CcDivSterW{} against \CcDivSterH), for laziness and for old age
(Fig.~\ref{fig:cultures}D). Almost every item is a failure of delivery, not of capacity: of provision, care and
restraint in men; of fidelity, fertility and work in women. This is Result~\ref{res:delivery} in the ethnographic
record. Capacity gets a person through the door, conduct keeps them inside, and the jury does not disband at the wedding.

\subsubsection{A caution: what a glance rewards is itself cultural} Of \CcBodyN{} societies with data, a plump or
moderately fat female body is the ideal in \CcBodyPlump\% and a slender one in \CcBodySlim\%~\cite{anderson1992}. Even
the most immediately legible trait is read against a local standard, so findings about ``attractiveness'' in one
society, including the speed-dating data used here, measure legibility and consensus, not a universal target.

\section{Today: believed, said, done, born}
\label{sec:today}

Two open data sets let us watch the links of the chain in living people: what each sex believes the other wants, what
people say they want, whom they choose after a brief meeting, and who ends up with children. Neither is ancestral. Both
show the structure the theory describes.

\subsection{Four-minute decisions}

\begin{table}[t]
\caption{What participants said mattered before the event (points out of 100) and what predicted a yes: change in the
probability of a yes, in percentage points, per standard deviation of the chooser's own rating of the partner, with
chooser fixed effects (95\% intervals are about $\pm2$ points); and the same using everyone else's rating of that
partner (``panel''). W = women choosing men; M = men choosing women.}
\label{tab:sd}
\small
\begin{tabularx}{\columnwidth}{@{}Lrrrrrr@{}}
\toprule
 & \multicolumn{2}{c}{Said} & \multicolumn{2}{c}{Did (own rating)} & \multicolumn{2}{c}{Did (panel)} \\
\cmidrule(lr){2-3}\cmidrule(lr){4-5}\cmidrule(l){6-7}
 & W & M & W & M & W & M \\
\midrule
Attractive & 17.9 & 27.0 & 15.3 & 21.5 & 11.7 & 16.0 \\
Shared interests & 12.7 & 10.9 & 10.3 & 7.8 & 2.9 & $-$0.5 \\
Fun & 17.3 & 17.6 & 6.8 & 10.1 & 3.5 & 5.1 \\
Intelligent & 20.9 & 19.4 & 4.1 & 1.2 & 3.0 & 2.5 \\
Sincere & 18.3 & 16.4 & 0.1 & 0.3 & $-$1.6 & $-$2.2 \\
Ambitious & 12.8 & 8.7 & 0.1 & $-$1.3 & $-$0.4 & 0.8 \\
\bottomrule
\end{tabularx}
\end{table}

\begin{figure*}[t]
\centering
\includegraphics[width=\textwidth]{figures/fig5_dating.pdf}
\caption{\textbf{Believed, said, done} (speed-dating data~\cite{fisman2006}; \SdN{} people, \SdRows{} dates).
\textbf{(A, B)}~For six traits: the weight the other sex believes it carries, the weight people say it carries (both as
points out of 100), and its share of the summed effects on a yes (Table~\ref{tab:sd}). \textbf{(C)}~Mean chance of a yes
that people expected from their partners (open) and received (filled), on dates with a recorded guess. \textbf{(D)}~Share of the variance in yes/no
decisions that belongs to the person judged, to the chooser, and to the pair or chance.}
\label{fig:dating}
\end{figure*}

\subsubsection{The sexes do much the same thing} A yes followed attractiveness above all (\SdRevAttM{} percentage points
per standard deviation for men choosing, \SdRevAttW{} for women), then shared interests and fun; intelligence earned
\SdRevIntW{} points from women and \SdRevIntM{} from men (Table~\ref{tab:sd}). Men respond more to looks (men minus women,
\SdRevDiffAtt{} points; interval \SdRevDiffCIAtt) and women more to intelligence (\SdRevDiffInt; \SdRevDiffCIInt), as the
original authors reported~\cite{fisman2006}, but the profiles have the same shape, as a meta-analysis of live
encounters also finds~\cite{eastwick2014}.

\subsubsection{Each sex reads the other better than the other reads itself} Participants also said what they thought
the other sex wanted. Men believed women give \DbMWBelAtt{} points to attractiveness; women said \DbMWSaidAtt; in the
event it carried \DbMWDidAtt\% of the weight in women's choices. Women believed men give \DbWMBelAtt{} points to
attractiveness; men said \DbWMSaidAtt; it carried \DbWMDidAtt\% (Fig.~\ref{fig:dating}A, B). For intelligence the order
reverses: believed \DbMWBelInt{} and \DbWMBelInt, said \DbMWSaidInt{} and \DbWMSaidInt, done \DbMWDidInt\% and
\DbWMDidInt\%. In both directions the sceptical belief of the other sex lies closer to behaviour than self-report does:
the mean absolute gap from what was done is \DbMWErrBel{} points for men's beliefs about women against \DbMWErrSaid{} for
women's own statements, and \DbWMErrBel{} against \DbWMErrSaid{} for women's beliefs about men. Stated preferences
describe the partner people want after long acquaintance; behaviour in an evening is bounded by what the evening shows
(Section~\ref{sec:legible}), and each sex has noticed this about the other.

\subsubsection{Men misread interest; nobody knows their own standing} After each date participants guessed the chance
that the partner would say yes. Men expected \DrMExp\% and received \DrMAct\%; women expected \DrWExp\% and received
\DrWAct\% (Fig.~\ref{fig:dating}C), the over-perception of interest by men that error-management theory
describes~\cite{haselton2000}. Neither sex could tell who in particular was interested: within persons, guesses
correlate only \DrMR{} and \DrWR{} with the partner's decision, but \DrMProj{} with one's own liking. People assume that
those they like, like them. Self-rated attractiveness exceeds the panel's rating by about a point in both sexes
(\DxMSelf{} against \DxMPanel{} for men; \DxWSelf{} against \DxWPanel{} for women).

\subsubsection{A yes belongs mostly to the pair} About a fifth of the variance in decisions belongs to the person judged
(\DvWTarget\% when women choose, \DvMTarget\% when men do), about a fifth to the chooser's general readiness to say yes,
and the rest to the particular pair or to chance (Fig.~\ref{fig:dating}D). Liking is barely reciprocal ($r=\DrLikeR$
across \DrPairs{} pairs), and mutual matches (\DrMatch\%) are no more common than independent choices would produce
(\DrMatchInd\%), which fits the finding that machine learning cannot predict who will desire whom from traits measured
beforehand~\cite{joel2017}. Consensus ``mate value'' is real but small. The market is nonetheless orderly:
people who were wanted more said yes less ($r=$~\DcMR{} for men and \DcWR{} for women; the most-wanted third of women said
yes to \DcWTop\% of partners and the least-wanted third to \DcWBot\%), and access was already unequal after four minutes,
more so for men (women said yes to \SdYesW\% of partners and men to \SdYesM\%; \SdPopZeroM\% of men and \SdPopZeroW\% of
women received no yes at all). Institutions matter here too: when women instead of men rotate between tables, the sex
difference in selectivity disappears~\cite{finkel2009}.

\subsection{Born: selection among 34,104 adults}

\begin{figure*}[t]
\centering
\includegraphics[width=\textwidth]{figures/fig6_episodes.pdf}
\caption{\textbf{Selection in the United States, split into three episodes} (General Social Survey 1972--2024; respondents
aged 45--79, born 1900--1979). \textbf{(A--C)}~Selection differentials on four traits, in standard deviations within sex
and decade of birth, decomposed by Result~\ref{res:episodes}, with 95\% bootstrap intervals (Table~\ref{tab:gss}).
\textbf{(D)}~Share of never-married respondents who are parents, by year of birth.}
\label{fig:gss}
\end{figure*}

\begin{table*}[t]
\caption{Selection differentials in the United States, decomposed into pairing (ever marrying), conversion (becoming a
parent) and number (family size among parents) by Result~\ref{res:episodes}. Respondents aged 45--79, born 1900--1979,
General Social Survey 1972--2024 ($n=\GeN$ with marital history); traits standardized and family size made relative within
sex and decade of birth; survey weights; [95\% bootstrap interval]. The three components sum to the total.}
\label{tab:gss}
\small
\begin{tabularx}{\textwidth}{@{}Llrlrlrlr@{}}
\toprule
Trait & Sex & \multicolumn{2}{c}{Pairing} & \multicolumn{2}{c}{Conversion} & \multicolumn{2}{c}{Number} & Total \\
\midrule
Personal income (log) & Men & +0.045 & [+0.037, +0.053] & +0.001 & [$-$0.009, +0.012] & $-$0.034 & [$-$0.051, $-$0.018] & +0.012 \\
 & Women & 0.000 & [$-$0.008, +0.007] & $-$0.037 & [$-$0.047, $-$0.027] & $-$0.079 & [$-$0.095, $-$0.063] & $-$0.116 \\
\addlinespace[2pt]
Occupational prestige & Men & +0.022 & [+0.017, +0.027] & $-$0.007 & [$-$0.013, +0.001] & $-$0.046 & [$-$0.057, $-$0.035] & $-$0.031 \\
 & Women & 0.000 & [$-$0.004, +0.005] & $-$0.020 & [$-$0.025, $-$0.014] & $-$0.062 & [$-$0.071, $-$0.052] & $-$0.082 \\
\addlinespace[2pt]
Vocabulary test & Men & +0.010 & [+0.001, +0.020] & $-$0.021 & [$-$0.033, $-$0.009] & $-$0.054 & [$-$0.072, $-$0.037] & $-$0.064 \\
 & Women & +0.004 & [$-$0.002, +0.010] & $-$0.032 & [$-$0.041, $-$0.023] & $-$0.090 & [$-$0.105, $-$0.074] & $-$0.117 \\
\addlinespace[2pt]
Years of schooling & Men & +0.014 & [+0.008, +0.020] & $-$0.017 & [$-$0.025, $-$0.011] & $-$0.081 & [$-$0.094, $-$0.069] & $-$0.085 \\
 & Women & $-$0.004 & [$-$0.009, 0.000] & $-$0.034 & [$-$0.040, $-$0.028] & $-$0.126 & [$-$0.137, $-$0.116] & $-$0.165 \\
\bottomrule
\end{tabularx}
\end{table*}

The General Social Survey has asked a probability sample of United States adults about schooling, work, income, marriage
and number of children since 1972~\cite{gss2024}. We take the \GssN{} respondents aged 45--79 and born 1900--1979
(\GssNW{} women, \GssNM{} men), standardize four traits within sex and decade of birth, and apply
Result~\ref{res:episodes} (Methods). Table~\ref{tab:gss} and Fig.~\ref{fig:gss} give the result.

\subsubsection{Men are selected at the door, women after it} For income, the pairing differential is \GeIncMPair{}
standard deviations in men (interval \GeIncMPairCI) and \GeIncWPair{} in women (\GeIncWPairCI). Conversion is
\GeIncMConv{} in men and \GeIncWConv{} in women (\GeIncWConvCI). Men who earn more are more often accepted as husbands
and gain nothing further; women who earn more are accepted equally often and then less often become mothers. The
pattern is general. For men all four traits are favoured at pairing (prestige \GePreMPair, schooling \GeEduMPair,
vocabulary \GeVocMPair; every interval excludes zero); for women none is.

\subsubsection{A floor for men, a ceiling for women} By fifths of income (Fig.~\ref{fig:main}, panel~5), childlessness
among men falls from \GeChlMOne\% in the lowest fifth to \GeChlMThree\% in the middle and \GeChlMFive\% in the highest,
and the share never married from \GeNevMOne\% to \GeNevMFive\%. The gradient is concentrated at the bottom (the fall
from lowest to middle exceeds that from middle to highest by \GeFloorM{} points; interval \GeFloorMCI). Among women
childlessness is flat across the lower three fifths (\GeChlWOne, \GeChlWTwo, \GeChlWThree\%) and rises to
\GeChlWFive\% at the top: a ceiling, the rise from middle to highest exceeding that from lowest to middle by
\GeFloorW{} points (\GeFloorWCI).

\subsubsection{Number is negative for everyone} Among parents, more schooling, a larger vocabulary, higher prestige and
higher income all go with fewer children, in both sexes, more steeply for women. For men's income the episodes cancel
(total \GssIncMTot): what the door rewards, the nursery takes back. For schooling and vocabulary the totals are negative
in both sexes. Timing carries much of it: each year's delay of the first birth goes with \GssDelayM\% (men) and
\GssDelayW\% (women) fewer children among parents, and age at first birth accounts statistically for \GssViaM\% and
\GssViaW\% of the schooling gradient.

\subsubsection{The gate is opening for women and narrowing for men} Among people born 1920--39, \GeGateMOld\% of
never-married men and \GeGateWOld\% of never-married women were parents; among those born 1960--79, \GeGateMNew\% and
\GeGateWNew\% (Fig.~\ref{fig:gss}D). Marriage no longer stands between women and motherhood, and still stands between
most men and fatherhood. Over the same cohorts the share of men who never married rose from \GeNeverMOld\% to
\GeNeverMNew\%, and the pairing differential on men's income tripled, from \GePairIncMOld{} to \GePairIncMNew{}
(\GePairIncMNewCI). For women it moved from \GePairIncWOld{} to \GePairIncWNew{} (\GePairIncWNewCI): in the youngest
cohort the door has begun, faintly, to read women's earnings as well. Childlessness now accounts for \GssShareNewM\% of the opportunity for selection among men and
\GssShareNewW\% among women, up from \GssShareOldM\% and \GssShareOldW\%.

\subsection{Mating behaviour, by sex}

\begin{figure*}[t]
\centering
\includegraphics[width=\textwidth]{figures/fig7_today.pdf}
\caption{\textbf{Mating in the United States, by sex} (General Social Survey). \textbf{(A)}~Reported opposite-sex partners
since age 18, adults aged 30--59, 1989--2024. \textbf{(B)}~Among married couples, the share in which the husband or the
wife has more years of schooling, by survey period. \textbf{(C)}~Childlessness at ages 45--79 by attendance at religious
services. \textbf{(D)}~Share aged 22--34 reporting no sexual partner in the past year; bars are 95\% intervals.}
\label{fig:today}
\end{figure*}

\subsubsection{Partner counts are skewed in both sexes, and are not the currency} Men aged 30--59 report a mean of
\GbMMean{} partners of the other sex since age 18 and a median of \GbMMed; women report \GbWMean{} and \GbWMed{}
(Fig.~\ref{fig:today}A). In a closed population the means must be equal, so the \GbMeanRatio-fold difference is
reporting and sampling, not behaviour, and counts should be compared within sex. Within each sex the distribution is
similarly unequal: the top fifth of men account for \GbMTopTwenty\% of all partners that men report and the top fifth of
women for \GbWTopTwenty\% of those women report (Gini \GbMGini{} and \GbWGini). The popular picture of a few men
monopolizing partners while women's experience is evenly spread finds no support here. Nor do partners turn into
children: men reporting one lifetime partner have \GssPartOneM{} children and men reporting twenty or more have
\GssPartManyM; for women the figures are \GssPartOneW{} and \GssPartManyW. Other differences are familiar. Among the
ever-married, \GbMAffair\% of men and \GbWAffair\% of women report an affair; \GbMPaid\% of men and \GbWPaid\% of women
report ever paying or being paid for sex; \GbMYearTwo\% of men and \GbWYearTwo\% of women aged 18--59 report two or more
partners in the past year.

\subsubsection{Serial marriage helps men a little} Among the currently married, men who have remarried after divorce
have \GbSerialMAbs\% more children than men married once (interval \GbSerialMCI), and remarried women
\GbSerialWAbs\% fewer (\GbSerialWCI), as found in another sample from the United States~\cite{jokela2010}. Remarriage is the form that
polygyny takes under monogamy, and it is modest.

\subsubsection{Who marries up has reversed} In 1972--82 the husband had more schooling than the wife in \GhHusbOld\% of
couples and the wife more in \GhWifeOld\%; in 2016--24 the figures are \GhHusbNew\% and \GhWifeNew\%
(Fig.~\ref{fig:today}B), the end of a pattern that has reversed across much of the world as women's schooling overtook
men's~\cite{esteve2016}. Customs of who marries whom follow the supply of partners. Spouses' years of schooling
correlate \GssSpOld{} in 1972--89 and \GssSpNew{} in 2008--24: assortment on a legible, certified trait is strong and
stable.

\subsubsection{A cultural contrast inside one country} Among men who attend religious services weekly,
\GrMWeeklyChl\% are childless at 45--79; among men who rarely or never attend, \GrMRareChl\%. For women the figures are
\GrWWeeklyChl\% and \GrWRareChl\% (Fig.~\ref{fig:today}C). Selection against schooling is markedly weaker among
weekly attenders (total differential \GrMWeeklyTot{} for men and \GrWWeeklyTot{} for women) than among those who rarely
attend (\GrMRareTot{} and \GrWRareTot). A congregation is a community in which people are watched, vouched for and
introduced, and the association is what the theory would expect; but attendance is measured after childbearing, and
having children brings people to services, so we draw no causal conclusion.

\subsubsection{The jury is being dismissed} Among adults aged 22--34 the share reporting no sexual partner in the past
year rose from \GssNoPOldM\% to \GssNoPNewM\% for men and from \GssNoPOldW\% to \GssNoPNewW\% for women between
1989--2002 and 2016--2024 (Fig.~\ref{fig:today}D), as others have found~\cite{ueda2020}. Over the same years the way
couples meet changed. Among heterosexual couples in the United States who met in 2017, 39\% met online, against 2\% in
1995; meeting through friends fell from 33\% to 20\% and through family from 15\% to 7\%; and almost nine in ten of those
who met online had no prior connection~\cite{rosenfeld2019}. On a dating site studied in four large cities, both sexes pursued partners
about 25\% more desirable than themselves~\cite{bruch2018}. In the terms of this paper, the people with years of observation are
leaving the room, and the instrument is shrinking toward a photograph and an evening. Section~\ref{sec:legible} says
what such an instrument reads. Whether this explains the narrowing door for men is prediction P9, not a finding.

\subsubsection{How to read these numbers} They are differences among phenotypes, not genes, in one country, and they are not
causal. Income is measured after most childbearing, and mothers' earnings fall after a first birth while fathers' do
not~\cite{kleven2019}, so part of the women's conversion differential runs from children to income. Men are known to
under-report children in surveys. Marriage is not the only form of pairing. The signs nevertheless hold unweighted, at
ages 50--74, and in surveys before and after 1990 (Appendix~B), and they match earlier work: positive selection on men's
income through marriage and childlessness and negative on women's in the United States and
Europe~\cite{hopcroft2006,hopcroft2021,fieder2007,weeden2006}; a positive and strengthening income gradient for Swedish men,
with very-low-income men increasingly childless~\cite{kolk2023}; and highest childlessness among the least-educated
Nordic men~\cite{jalovaara2019}. Genetic studies find that variants associated with more schooling are now selected
against~\cite{beauchamp2016,kong2017,hughjones2022}, as the negative number differential would lead one to expect.

\subsubsection{What the present shows} Contraception, schooling and later marriage have changed the later episodes beyond
recognition: where status once brought more surviving children~\cite{vonrueden2016}, it now brings fewer. The present is
therefore no guide to the \emph{direction} of past selection. It is a guide to its \emph{structure}. The door still
filters men on standing and earning, as in the ethnographic record; what follows still weighs most on women; and the
episodes are decided by different people, at different ages, on different evidence.

\section{Drivers of brain and culture}
\label{sec:drivers}

Equations~\eqref{eq:identity} and~\eqref{eq:optimum} turn the question ``what drove the human brain?'' into a list with
units. A candidate driver must have moved accuracy $\rho$ (through the number of readings $k$, their independence,
their fairness or their breadth $K$), alignment, opportunity $I$, the direct payoff $\beta_{\mathrm{own}}$, or the cost
$\kappa$. Table~\ref{tab:drivers} sorts the main candidates this way and grades the evidence.

\begin{table*}[t]
\caption{Candidate drivers, the term each moves, and the strength of the evidence. ``Strong'' means well documented in
living or recent societies or in dated material; whether the practice had this effect in the deep past is an inference in
every row.}
\label{tab:drivers}
\small
\begin{tabularx}{\textwidth}{@{}P{0.16\textwidth}LP{0.13\textwidth}P{0.25\textwidth}@{}}
\toprule
Driver or behaviour & What it made readable, or changed & Term moved & Evidence \\
\midrule
Shared care of children & A mother's success depends on recruited helpers, who choose whom to help; partners are judged on reliability & $\beta_{\mathrm{social}}\uparrow$, $k\uparrow$ & Strong for help and survival~\cite{sear2008,hrdy2009,kramer2010}; link to cognition comparative~\cite{burkart2014} \\
Skill-intensive foraging and hunting & Long learning curves, returns visible to the camp, meat shared in public & $\beta_{\mathrm{own}}\uparrow$, $K\uparrow$ & Strong~\cite{kaplan2000,koster2020,gurven2006,smith2004} \\
Fire and cooking & Cheaper energy for brain tissue; lit evenings add hours of talk & $\kappa\downarrow$, $k\uparrow$ & Dates strong ($\sim$1.0 and 0.78 Ma)~\cite{berna2012,zohar2022}; energetics experimental~\cite{carmody2011}; evening talk~\cite{wiessner2014} \\
Language and storytelling & Verbal skill displayed nightly; reputations travel beyond eyewitnesses & $k\uparrow$, $K\uparrow$, but errors shared & Moderate~\cite{smith2017,wiessner2014,dunbar2004} \\
Play, contests and games & Rule-bound, safe, repeatable comparisons among age-mates & fairness $\uparrow$, $k\uparrow$ & Strong in ethnography~\cite{lewlevy2017,deaner2013}; board games attested only from $\sim$5,000 years ago~\cite{crist2016} \\
Arranged marriage, bride service and the wait & A jury of in-laws; an audition of years; unstaged observation of conduct & $k\uparrow$, delivery read, floor on men & Strong: this paper (Section~\ref{sec:cultures}); $\geq$50,000 years by phylogeny~\cite{walker2011} \\
Levelling coalitions and divorce for cruelty & Bullies sanctioned; deference flows to skill and generosity & aim: from force to competence & Moderate~\cite{boehm1999,henrich2001,wrangham2019}; Fig.~\ref{fig:cultures}D \\
Stored wealth and rank & Extra wives and standing go to inheritance as well as skill & circumstance $\uparrow$, accuracy $\downarrow$ & This paper (Fig.~\ref{fig:cultures}B); effect on selection untested \\
Division of labour, credentials, contraception, online meeting & One specialty per person; outcomes certified by institutions; family size chosen; the jury dismissed & $K\downarrow$, $k\downarrow$, number episode reversed & This paper (Section~\ref{sec:today}); otherwise speculative \\
\bottomrule
\end{tabularx}
\end{table*}

\emph{Culture is the instrument, not only the product.} Games, stories, contests and courtship rituals are usually
treated as things a large brain made possible. On the present account they are also measuring devices that made
selection for a large brain possible, by raising $\rho$ at the age of decision. Instrument and capacity can then
ratchet: better minds invent richer performances, and richer performances read minds better.

\emph{The same driver can act on both sides of the ratio.} Fire lowered the cost of tissue and, by lighting the evening,
lengthened the daily window for performance~\cite{wiessner2014}. The steps in the fossil curve near \FosBreakA{} and
\FosBreakB{} million years ago predate the earliest secure fire, so fire cannot explain them; it may bear on the later
rise.

\emph{The optimum can stall, and volume and capacity can part.} Brain volume did not rise within \Hs{} over 300,000
years and appears to have fallen in the last 12,000 (Section~\ref{sec:puzzle}), the period in which stored wealth, a
division of labour and written records changed who is judged, by whom and on what. Yet an ancient-DNA study of 15,836
West Eurasians reports that, over the same ten millennia, combinations of alleles that today predict measures of
cognitive performance rose by about one standard deviation on the scale of modern variation; the authors caution that
such effects were measured in industrialized societies and that their relation to what was adaptive in the past is
unclear~\cite{akbari2026}. If both findings hold, capacity and volume have come apart, one more reason to treat volume as
a blunt proxy, and the pace of that change, far above the Pleistocene average, is what Result~\ref{res:floor} says weak
filters can deliver when the optimum moves. In the United States today pairing still runs in the ancestral direction
for men, the later episodes run against it for both sexes, and the totals are negative (Table~\ref{tab:gss}). None of
this is a forecast. Differentials of a tenth of a standard deviation on phenotypes, with modest heritability and
cultural transmission of schooling in the mix, imply genetic change that is small and easily reversed by a change in
when people have children.

\begin{table*}[t]
\noindent\colorbox{band}{\parbox{\dimexpr\textwidth-2\fboxsep}{\small
{\sffamily\bfseries Box 1. The same problem in artificial minds.}\;
A machine-learning model is chosen, funded and released on its performance in tests, long before the outcomes that
matter, its usefulness and safety over years of deployment, can be known. Evaluation is performance before outcomes,
and the eight results carry over term by term.\par\vspace{3pt}
\footnotesize
\begin{tabularx}{\linewidth}{@{}P{0.155\linewidth}P{0.30\linewidth}L@{}}
\toprule
Result & In human evolution & In the selection of machine-learning systems \\
\midrule
\ref{res:identity}. Accuracy $\times$ alignment $\times$ opportunity & Response is variation, times the accuracy of the early reading, times the fidelity of decisions to it, times reproductive inequality & Progress per generation of models is variation among candidates, times the validity of the evaluations, times the fidelity of decisions to them, times the strength of selection. Better evaluation substitutes for harsher selection \\
\ref{res:jury}. Juries, and their ceiling & Many observers and occasions average out error; gossip shares its errors & Independent evaluations average out; a test set that has leaked into training or been tuned against for years is gossip, capped at $r_1/(r_1+c(1-r_1))$. Fresh held-out tests lift the cap~\cite{dwork2015,recht2019} \\
\ref{res:fair}. Fair contests & Same task, same day, among age-mates: inheritance and luck are stripped out & Identical prompts, tools and budgets: a small standardized test can say more about a model than a large uncontrolled record \\
\ref{res:delivery}. Display and delivery & Skill shows in an audition; reliability only in unwatched conduct, and is never proved & A test the system can recognize measures what it can do, not what it will do. Honesty and restraint can be read only on occasions it cannot tell from ordinary use, or must be enforced from outside; after $m$ clean occasions the failure rate is bounded only by $3/m$ \\
\ref{res:breadth}. Breadth & Judged as hunter, maker, storyteller and parent, a person is selected on what those share & Averaging $K$ unrelated tasks aims selection $K\sqrt{V_G/V_S}$ times more sharply at what they share: broad evaluation breeds general systems; one leaderboard breeds specialists and, once it is the target, ceases to measure~\cite{campbell1979} \\
\ref{res:floor}. A floor is enough & Passing over the unready suffices; no despot is needed & Discarding the worst candidates is enough; winner-take-all contests add conflict, not much speed \\
\ref{res:optimum}. Reading $\times$ energy & Favoured brain size is legibility times payoff, over metabolic cost & Capability worth building is what can be read, times what it earns, over what it costs: cheaper computation is worth most where evaluation is good \\
\ref{res:episodes}. Episodes & Pairing, conversion and number have different judges and different evidence & Release and adoption are different episodes, judged by different people on different evidence \\
\bottomrule
\end{tabularx}\par\vspace{3pt}
\small Human evolution ran this experiment for three million years. Its lesson for those who build minds is that the
instrument sets the target: we will get the kind of intelligence our evaluations can read, and the kind of character we
can observe unannounced or are prepared to enforce.}}
\end{table*}

\section{Predictions}
\label{sec:predictions}

The account would be damaged by any of the following failing.

\begin{enumerate}[label=\textbf{P\arabic*.},leftmargin=2.2em]
\item \textbf{Jury law in the field.} In foraging camps, agreement about a person's skill, character and intelligence
rises with the number of raters as Eq.~\eqref{eq:sb} predicts, and multi-rater reputations in adolescence predict adult
returns and reproductive success better than any single early outcome.
\item \textbf{Early performance predicts late outcomes.} Rankings in children's and adolescents' play and contests predict
adult production and surviving children 15 or more years later, after allowing for parents' wealth. A validity
below 0.1 in a well-powered longitudinal sample refutes the mechanism.
\item \textbf{Fair contests are more heritable.} Within one population, rankings from rule-bound contests among age-mates
show higher parent--offspring resemblance than rankings from accumulated outcomes (Result~\ref{res:fair}).
\item \textbf{Pairing for men, later episodes for women.} In natural-fertility populations the pairing term of
Result~\ref{res:episodes} dominates the selection differential on men's skill and standing, and conversion and number
(through child survival and birth intervals) dominate for women's.
\item \textbf{A floor and a queue.} The relation between men's standing and reproductive success is concave, steepest
between the unmarried and the married; and within a society, lower-ranked men marry later, not only less.
\item \textbf{Breadth and the general factor.} Where partners are judged across more domains, choice loads more
heavily on general ability; where occupations are specialized and certified, it loads less
(Result~\ref{res:breadth}).
\item \textbf{Who judges sets what is selected.} Where parents choose, spouses resemble each other more in family
standing and reliability and less in looks than where individuals choose; and within a society, arranged and
self-chosen couples differ in the same way.
\item \textbf{Legibility bounds preference.} In any brief-encounter data, the weight a trait carries in choice tracks
how well one encounter reads it (Fig.~\ref{fig:early}B); lengthening encounters or adding observers shifts weight toward
intelligence and sincerity.
\item \textbf{The dismissed jury.} As meeting moves from friends and family to strangers online, selection at the door
shifts toward what a profile shows, and men with poor profiles but good reputations lose most. Couples introduced by
people who know both are more alike in rated character than couples who met as strangers.
\item \textbf{Steps follow instruments or energy.} Each sustained rise in endocranial volume follows, and does not
precede, evidence for a fall in the cost of tissue or a rise in legibility (shared care, public provisioning, lit
evenings, symbolic display). A step with neither counts against the account.
\item \textbf{Pulses, not a steady push.} In ancient genomes, change in cognition-associated variants is episodic, and its
pulses coincide with changes in who is compared with whom and at what stakes (density, exchange, record-keeping), not
only with climate. The first large study finds such a pulse in the last ten millennia~\cite{akbari2026}; its timing within
that span is the test.
\item \textbf{The same law in machines.} Across families of models, gains on a benchmark whose test items are shared
or reused saturate at the ceiling of Result~\ref{res:jury}; and the generality of the systems a field produces rises
with the breadth of the evaluations it selects on (Box~1).
\end{enumerate}

\noindent A direct experimental test is also possible: randomize observers to one, four or sixteen viewings of either a
standardized performance or a record of outcomes for the same individuals, and score each condition against a criterion
measured later. Results~\ref{res:jury} and~\ref{res:fair} predict the size of every cell from two parameters.

\section{Limits and objections}
\label{sec:limits}

\subsubsection{No ancestral parameter is estimated} We do not know $\rho$, $I$ or $\kappa$ for any Pleistocene population.
The fossil record fixes the rate and the pattern; the ethnographic record shows that the institutions exist and what
they reward; the two modern data sets measure the links in living people. The inference from these to deep time is an
argument, and Section~\ref{sec:predictions} is the list of ways to break it.

\subsubsection{Recorded societies are not ancestors} They were described in the last two centuries, mostly by outsiders
and mostly by men, they live beside farmers and states, and they are related to one another. The codes compress customs
into categories. We rely on large samples and on counts, and treat every such finding as a tendency.

\subsubsection{Volume is not capacity} Section~\ref{sec:puzzle} gives the caveat. A theory about minds tested against
skulls has slack in it.

\subsubsection{The modern samples are unusual} Graduate students at one university in 2002--2004 and adults in one
country are not humanity. The jury law is a statistical necessity whenever errors are independent, so we expect it to
travel; the particular reliabilities, and the sizes and even signs of the modern selection differentials, need not.

\subsubsection{Relation to ecological, social and cultural accounts} Those accounts explain why capacity pays
($\beta_{\mathrm{own}}$, $\beta_{\mathrm{social}}$) and what it costs ($\kappa$). Ours adds the term that converts
social payoffs into selection ($\rho$ at the age of decision) and so predicts things they do not: that selection aims at
general ability in proportion to the breadth of what is judged, that juries, auditions and waits matter, that fairness
of contests matters, and that a preference cannot outrun its instrument. If those predictions fail, the older accounts
stand without amendment.

\subsubsection{Sexual selection by another name} Partly. But the judges here are as often parents, in-laws,
helpers and allies as mates, the currency is as often help and tolerance as copulation, and the filter is a floor.
West-Eberhard's broader term, social selection~\cite{westeberhard1983,nesse2007}, fits better.

\subsubsection{Faking} A display that can be faked cheaply stops being believed. Repetition, publicity and common rules
are what make faking expensive: it is hard to appear to be a good hunter, storyteller or mother every day for years in
front of the same thirty people. Result~\ref{res:delivery} says where faking remains possible: in conduct, on occasions
known to be tests.

\subsubsection{The analogy with machines} Box~1 transfers mathematics, not biology. Models are not organisms and
developers are not mates. The results apply because they are statements about selection on scores, whatever is being
selected.

\section{Conclusion}

A capacity that pays late can be selected early only if it can be read early. Humans built the instruments for reading
it: play, contests, shared work, stories by the fire, juries of kin, auditions for in-laws, and a wait. The fossil
record leaves ample room for a mechanism of this kind, since the rate of brain expansion was a hundredth of what even a
mild filter delivers, and it has the right shape, steps and stasis, for a trait tracking an optimum that such
instruments would move. We have measured the law by which observers converge on a verdict and found that reading a mind
takes a camp, not a suitor; that people choose on what they can read; and that the world's marriage customs are
arranged accordingly, with a jury that decides, a groom who pays, serves and waits, and a marriage that lasts only as
long as conduct does. Selection divides into episodes that are judged by different people on different evidence, and
they have weighed differently on men and on women, in the ethnographic record and in the United States today. Judging
people across many arenas aims selection at general ability.

The oldest question about human intelligence is what it was for. The data and results here suggest an addition: it
mattered just as much that intelligence could be \emph{seen}, early enough to count, by the people whose decisions
shaped the next generation. We did not have to wait for the end of a life to know what was in it. We are now building
minds of another kind, and choosing among them, again, by performance before outcomes. What we get will depend, as it
did the first time, on how well we learn to read.

\section*{Methods}
\addcontentsline{toc}{section}{Methods}
\small

\subsubsection{Fossil crania} The file is the public compilation of DeSilva et~al.~\cite{desilva2021}: 987 rows, of which two
(a Miocene ape) were dropped, leaving \FosN{} hominin crania with volume, date and source. Species labels were grouped as
in Table~\ref{tab:fossils}. Rates use the \FosNRate{} crania dated between 3.2 million and 11,700 years ago, excluding
\textit{H.~naledi} and \textit{H.~floresiensis}. The specimen-level slope is an ordinary least-squares regression of log
volume on age with a 2,000-replicate bootstrap; the binned slope regresses the means of \FosNBins{} 100,000-year bins. The
within-bin standard deviation pools bins with at least three crania. Rates in standard deviations per generation divide
the slope by the standard deviation and by 40,000 generations per million years. Continuous piecewise-linear models with
one, two and three breaks were fitted by grid search (steps of 0.05 or 0.1 million years; adjacent breaks at least 0.2
apart) and compared by AICc and BIC; intervals for the two breaks come from 300 bootstrap refits. As checks, the file
reproduces the published Pleistocene \Hs{} mean and standard deviation~\cite{desilva2023}, and \FosDuMatch{} of the
\FosDuN{} specimens of Du et~al.~\cite{du2018} match by volume and date ($r=\FosDuR$).

\subsubsection{Speed dating} The public file of the Columbia experiment~\cite{fisman2006} has \SdRows{} rows, one per
person per date. Ratings are on a 1--10 scale. \emph{Agreement}: ratings were centred on each rater's mean; for each
target with at least $2k$ raters, two disjoint random sets of $k$ raters were averaged and the two averages
correlated across targets; values are means of 400 random draws. \emph{Said} and \emph{believed}: importance of six
attributes to oneself, and the importance one thinks the other sex gives them, stated before the event and rescaled to
sum to 100 per person (four sessions used a 1--10 scale). \emph{Did}: a linear probability model of the yes/no decision
on the six ratings, standardized over all dates, with a fixed effect for each chooser, fitted separately for women and
men; intervals from 600 bootstrap resamples of choosers; shares are coefficients divided by their sum. The ``panel''
version replaces a chooser's own ratings by the mean rating the partner received from all other choosers. Expected
interest is the stated probability (1--10) that the partner would say yes. Variance shares are one-way
analysis-of-variance estimates for the person judged, after removing chooser means, and for the chooser, after removing
target means. The college entrance-test measure is the median SAT score of the participant's undergraduate institution,
as coded by the original authors.

\subsubsection{Cross-cultural codes} We used the D-PLACE release~\cite{kirby2016} of the Ethnographic
Atlas~\cite{murdock1967} (variables EA006, EA009, EA042), the Standard Cross-Cultural Sample~\cite{murdock1969} (marriage
arrangements for males and females, SCCS739--740~\cite{broude1983}; relative age at first marriage,
SCCS612~\cite{whyte1978}; polygyny, SCCS866--871~\cite{white1986}; causes of divorce, SCCS1136--1175~\cite{betzig1989};
consent, SCCS971--972~\cite{frayser1985}; ideal female body type, SCCS1248~\cite{anderson1992}) and Binford's
hunter-gatherer data~\cite{binford2001} (B019--B021). Missing codes were dropped variable by variable, so sample sizes
differ and are given in the text. Foragers are societies whose dominant subsistence is gathering, hunting or fishing.
For marriage arrangements, codes 5--6 are ``parents choose'', 2--3 ``individual chooses with approval'', 4 ``either''.
For divorce, each cause is coded by the spouse at fault. Ages at first marriage in Binford's compilation are ethnographers'
estimates, often rounded, so levels are less secure than the difference between the sexes. The intervals for the forager
age gap and for its correlation with polygyny are bootstraps over groups and ignore their relatedness.

\subsubsection{General Social Survey} We read the 1972--2024 cumulative file~\cite{gss2024} as distributed in the
\texttt{gssr} package~\cite{healygssr} (\GssNFile{} respondents). The analysis sample is respondents aged 45--79 and born
1900--1979 with known sex and number of children (\GssN). Traits: years of schooling, the ten-word vocabulary test,
occupational prestige (2010 scale) and the logarithm of personal income in constant dollars (ages 45--64, positive
incomes). Each trait was standardized, and number of children divided by its mean, within sex and decade of birth, using
the post-stratification weight. Number of children is top-coded at eight. Pairing is ever having married. Intervals are
percentiles of 1,000 Poisson bootstrap replicates, with the standardization repeated inside each replicate. Income fifths are weighted fifths within sex and decade of birth.
Lifetime partner counts are opposite-sex partners since age 18 (adults aged 30--59, surveys from 1989). Remarried
respondents are the currently married who report a past divorce or separation. Attendance groups are weekly or nearly
(codes 6--8), sometimes (2--5) and rarely or never (0--1).

\subsubsection{Simulations} Each formal result was checked by Monte Carlo with up to $2\times10^6$ individuals: the
identity of Result~\ref{res:identity} under four rules (a floor, a doubled top fifth, an exponential and a rank rule;
largest error \ThIdentErr), the jury law (largest error \ThJuryErr), the covariances and reliabilities of
Result~\ref{res:delivery} (\ThDelivErr), the breadth ratio, and the floor intensities. All random seeds are fixed.

\subsubsection{What is not computed here} Statements about particular foraging societies, life histories, genetics and
other populations are quoted from the cited sources, with the figures as published; we did not reanalyse those data.

\section*{Data and code}
All analyses run from the accompanying package: \texttt{code/} (Python), \texttt{data/} (the open files as downloaded, or extracts of the two largest made by the scripts provided),
\texttt{results/} (JSON and CSV) and \texttt{figures/}. Every number in the text that derives from data is written into
the manuscript by \texttt{code/make\_numbers.py}. Sources: the endocranial-volume tables published with
DeSilva et~al.~\cite{desilva2021} and Du et~al.~\cite{du2018}, and the speed-dating file and key released by Fisman and
Iyengar~\cite{fisman2006}, each obtained from a public mirror and checked against the publications as described above
(addresses and checksums are in the package); D-PLACE~\cite{kirby2016} from its public repository; the General Social
Survey from NORC~\cite{gss2024} through the \texttt{gssr} package~\cite{healygssr}.

\section*{Acknowledgments and AI-assistance disclosure}
The idea and its direction are the author's. Literature search, mathematics, code, data analysis, figures and drafting
were carried out with Claude (Anthropic) under the author's direction; earlier drafts also used ChatGPT (OpenAI). No
figure contains generated imagery: Figures~\ref{fig:main}, \ref{fig:early}B, \ref{fig:cultures}, \ref{fig:dating},
\ref{fig:gss} and~\ref{fig:today} plot measured data, Figure~\ref{fig:early}A plots published values, and
Figure~\ref{fig:model} plots the model. Quantities taken from the literature were checked against the sources where
these could be accessed; the author is responsible for the content and will verify each citation against the original
before any journal submission. The work has not been peer reviewed.

\appendix
\section*{Appendix A: Proofs}
\addcontentsline{toc}{section}{Appendix A}

\noindent\textbf{Result 1.} By the law of total covariance,
$\Cov(a,w)=\E[\Cov(a,w\mid Y)]+\Cov(\E[a\mid Y],\E[w\mid Y])$. The first term is zero because $w$ and $a$ are
independent given $Y$. With $\E[a\mid Y]=\bar a+b\,(Y-\bar Y)$ and $b=\Cov(a,Y)/\Var(Y)$, the second is
$b\,\Cov(Y,\E[w\mid Y])=b\,\Cov(Y,w)$. Hence $\Cov(a,w)=\Cov(a,Y)\Cov(Y,w)/\Var(Y)=\sigma_A\rho_{aY}\rho_{Yw}\sigma_w$,
and $\sigma_w=\sqrt{I}$ since $\bar w=1$. If $Y$ depends on $a$ only through $z$, then
$\rho_{aY}=\rho_{az}\rho_{zY}=h\sqrt{R}$. No assumption is made about the form of the rule linking $Y$ to $w$. If
capacity also raises fitness directly, the direct covariance adds to the right-hand side.

\noindent\textbf{Result 2.} Write a reading as $Y_j=t+\varepsilon_j$ with $\Var(t)=r_1$, $\Var(\varepsilon_j)=1-r_1$ and
$\mathrm{Corr}(\varepsilon_i,\varepsilon_j)=c$. The mean of $k$ readings has variance
$r_1+(1-r_1)\,[1/k+(1-1/k)\,c]$, of which $r_1$ is true variance; the ratio is the reliability. Setting $c=0$ gives
Eq.~\eqref{eq:sb}; letting $k\to\infty$ gives the ceiling.

\noindent\textbf{Result 3.} $\Cov(a,Y)=\sigma_A^2$ and $\Var(Y)=\sigma_A^2+\sigma_u^2+\sigma_\varepsilon^2$ for independent
components. Cue 1 has the higher $\rho_{aY}$ exactly when $\sigma_{u,1}^2+\sigma_{\varepsilon,1}^2<\sigma_{u,2}^2+\sigma_{\varepsilon,2}^2$.

\noindent\textbf{Result 4.} (i)~With $b$ and $p$ independent, $\Cov(b,pb)=\E[p]\,\E[b^2]-\E[p]\,\E[b]^2=\bar p\,V_B$, and
symmetrically for $p$. (ii)~On a recognized test each person performs at capacity, so the $dn$ staged occasions are
readings $b+\varepsilon$ with error variance $\sigma^2$, whose mean has reliability $V_B/(V_B+\sigma^2/(dn))$; on a
recognized test conduct carries no information about $p$, so only the $(1-d)n$ unstaged occasions, with error variance
$\nu$, inform about it. (iii)~If the default rate is $\theta$, the chance of $m$ clean occasions is $(1-\theta)^m$,
which exceeds 0.05 for all $\theta<1-0.05^{1/m}\approx3/m$.

\noindent\textbf{Result 5.} With $Y=K^{-1}\sum_k x_k$, $\Cov(Y,g)=V_G$ and $\Cov(Y,s_k)=V_S/K$. For any rule based on
$Y$ under joint normality, the selection differential on a component $v$ is $\Cov(v,Y)\,S_Y/\Var(Y)$; dividing by the
component's standard deviation and taking the ratio gives $K\sqrt{V_G/V_S}$. The validity follows from
$\Var(Y)=V_G+(V_S+V_\varepsilon)/K$.

\noindent\textbf{Result 6.} For a standard normal score truncated below at $z_q=\Phi^{-1}(q)$, the mean of the retained
part is $\phi(z_q)/(1-q)$. The response in a trait with $\rho_{zY}=\sqrt{R}$ is $h^2\sqrt{R}\,i(q)$ when both sexes are
selected and half that when one is~\cite{falconer1996}.

\noindent\textbf{Result 7.} $\partial\ln w/\partial B=a\beta/B-\kappa=0$ gives $B^*=a\beta/\kappa$; the second derivative
is negative. More generally, for benefit $\beta f(B)$ with $f'>0>f''$ and cost $\kappa\,c(B)$ with $c'>0$, $c''\ge0$,
the implicit-function theorem gives $\partial B^*/\partial\beta>0$ and $\partial B^*/\partial\kappa<0$.

\noindent\textbf{Result 8.} With $\E z=0$ and $\E c=1$, $S=\E[zc]=p\,\E[zc\mid c>0]=p\,\{\Cov(z,c\mid c>0)+\E[z\mid c>0]\,
\E[c\mid c>0]\}$, and $p\,\E[c\mid c>0]=1$. Adding and subtracting $\E[z\mid\text{ever paired}]$ splits access into pairing
and conversion. For the opportunity, $I=\E[c^2]-1=p\,\E[c^2\mid c>0]-1$, with
$\E[c^2\mid c>0]=(1+I_{\mathrm{parents}})/p^2$.

\section*{Appendix B: Robustness of the decomposition}
\addcontentsline{toc}{section}{Appendix B}
Table~\ref{tab:rob} repeats the decomposition of Table~\ref{tab:gss} without survey weights, in a narrower age range, in
two survey periods and in two birth cohorts. Every sign in the main specification is preserved in every row.

\begin{table}[h]
\caption{Access and number differentials (standard deviations) under alternative specifications. Access is pairing plus
conversion. Inc = personal income; Pre = occupational prestige; Sch = years of schooling; Voc = vocabulary. M = men,
W = women.}
\label{tab:rob}
\scriptsize
\setlength{\tabcolsep}{2.2pt}
\begin{tabularx}{\columnwidth}{@{}Lrrrrrrrr@{}}
\toprule
 & \multicolumn{4}{c}{Access} & \multicolumn{4}{c}{Number} \\
\cmidrule(lr){2-5}\cmidrule(l){6-9}
 & Inc M & Inc W & Pre M & Pre W & Sch M & Sch W & Voc M & Voc W \\
\midrule
Main specification & +0.046 & $-$0.037 & +0.015 & $-$0.019 & $-$0.081 & $-$0.126 & $-$0.054 & $-$0.090 \\
Unweighted & +0.048 & $-$0.053 & +0.013 & $-$0.029 & $-$0.082 & $-$0.122 & $-$0.064 & $-$0.089 \\
Ages 50--74 only & +0.044 & $-$0.033 & +0.016 & $-$0.019 & $-$0.087 & $-$0.122 & $-$0.059 & $-$0.096 \\
Surveys 1972--1990 & +0.041 & $-$0.038 & +0.018 & $-$0.018 & $-$0.104 & $-$0.119 & $-$0.094 & $-$0.092 \\
Surveys 1991--2024 & +0.049 & $-$0.037 & +0.014 & $-$0.020 & $-$0.071 & $-$0.129 & $-$0.037 & $-$0.088 \\
Born 1920--1939 & +0.041 & $-$0.030 & +0.011 & $-$0.017 & $-$0.070 & $-$0.116 & $-$0.072 & $-$0.083 \\
Born 1940--1959 & +0.039 & $-$0.039 & +0.007 & $-$0.024 & $-$0.083 & $-$0.134 & $-$0.028 & $-$0.099 \\
\bottomrule
\end{tabularx}
\end{table}

\normalsize

\begingroup\footnotesize\setlength{\bibsep}{1.5pt plus 0.5pt}
\begin{thebibliography}{118}

\bibitem[{DeSilva} et~al.(2021){DeSilva}, {Traniello}, {Claxton}, and
  {Fannin}]{desilva2021}
J.~M. {DeSilva}, J.~F.~A. {Traniello}, A.~G. {Claxton}, and L.~D. {Fannin}.
\newblock {When and why did human brains decrease in size? A new change-point
  analysis and insights from brain evolution in ants}.
\newblock \emph{Front. Ecol. Evol.}, 9:\penalty0 742639, 2021.
\newblock \doi{10.3389/fevo.2021.742639}.

\bibitem[{DeSilva} et~al.(2023){DeSilva}, {Fannin}, {Cheney}, {Claxton},
  {Ilieş}, {Kittelberger}, {Stibel}, and {Traniello}]{desilva2023}
J.~{DeSilva}, L.~{Fannin}, I.~{Cheney}, A.~{Claxton}, I.~{Ilieş},
  J.~{Kittelberger}, J.~{Stibel}, and J.~{Traniello}.
\newblock {Human brains have shrunk: the questions are when and why}.
\newblock \emph{Front. Ecol. Evol.}, 11:\penalty0 1191274, 2023.
\newblock \doi{10.3389/fevo.2023.1191274}.

\bibitem[{Fisman} et~al.(2006){Fisman}, {Iyengar}, {Kamenica}, and
  {Simonson}]{fisman2006}
R.~{Fisman}, S.~S. {Iyengar}, E.~{Kamenica}, and I.~{Simonson}.
\newblock {Gender differences in mate selection: Evidence from a speed dating
  experiment}.
\newblock \emph{The Quarterly Journal of Economics}, 121\penalty0 (2):\penalty0
  673--697, 2006.
\newblock \doi{10.1162/qjec.2006.121.2.673}.

\bibitem[{Broude} and {Greene}(1983)]{broude1983}
G.~J. {Broude} and S.~J. {Greene}.
\newblock {Cross-cultural codes on husband-wife relationships}.
\newblock \emph{Ethnology}, 22\penalty0 (3):\penalty0 263--280, 1983.
\newblock \doi{10.2307/3773467}.

\bibitem[{Kirby} et~al.(2016)]{kirby2016}
K.~R. {Kirby}, R.~D. {Gray}, S.~J. {Greenhill}, F.~M. {Jordan}, S.~{Gomes-Ng}, H.-J. {Bibiko}, et~al.
\newblock {D-PLACE: A global database of cultural, linguistic and environmental diversity}.
\newblock \emph{PLoS ONE}, 11\penalty0 (7):\penalty0 e0158391, 2016.
\newblock \doi{10.1371/journal.pone.0158391}.

\bibitem[{Binford}(2001)]{binford2001}
L.~R. {Binford}.
\newblock \emph{{Constructing Frames of Reference: An Analytical Method for Archaeological Theory Building Using
  Ethnographic and Environmental Data Sets}}.
\newblock University of California Press, Berkeley, 2001.

\bibitem[{Davern} et~al.(2025){Davern}, {Bautista}, {Freese}, {Herd}, and
  {Morgan}]{gss2024}
M.~{Davern}, R.~{Bautista}, J.~{Freese}, P.~{Herd}, and S.~L. {Morgan}.
\newblock {General Social Survey 1972--2024}.
\newblock Machine-readable cumulative data file. NORC at the University of
  Chicago, Chicago. \url{https://gss.norc.org}, 2025.

\bibitem[{Kuzawa} et~al.(2014){Kuzawa}, {Chugani}, {Grossman}, {Lipovich},
  {Muzik}, {Hof}, et~al.]{kuzawa2014}
C.~W. {Kuzawa}, H.~T. {Chugani}, L.~I. {Grossman}, L.~{Lipovich}, O.~{Muzik},
  P.~R. {Hof}, et~al.
\newblock {Metabolic costs and evolutionary implications of human brain
  development}.
\newblock \emph{Proc. Natl. Acad. Sci. USA}, 111\penalty0 (36):\penalty0
  13010--13015, 2014.
\newblock \doi{10.1073/pnas.1323099111}.

\bibitem[{Aiello} and {Wheeler}(1995)]{aiello1995}
L.~C. {Aiello} and P.~{Wheeler}.
\newblock {The expensive-tissue hypothesis: The brain and the digestive system in human and primate evolution}.
\newblock \emph{Curr. Anthropol.}, 36\penalty0 (2):\penalty0 199--221, 1995.
\newblock \doi{10.1086/204350}.

\bibitem[{Carmody} et~al.(2011){Carmody}, {Weintraub}, and
  {Wrangham}]{carmody2011}
R.~N. {Carmody}, G.~S. {Weintraub}, and R.~W. {Wrangham}.
\newblock {Energetic consequences of thermal and nonthermal food processing}.
\newblock \emph{Proc. Natl. Acad. Sci. USA}, 108\penalty0 (48):\penalty0
  19199--19203, 2011.
\newblock \doi{10.1073/pnas.1112128108}.

\bibitem[{Kaplan} et~al.(2000){Kaplan}, {Hill}, {Lancaster}, and
  {Hurtado}]{kaplan2000}
H.~{Kaplan}, K.~{Hill}, J.~{Lancaster}, and A.~M. {Hurtado}.
\newblock {A theory of human life history evolution: Diet, intelligence, and
  longevity}.
\newblock \emph{Evol. Anthropol.}, 9\penalty0 (4):\penalty0 156--185, 2000.
\newblock \doi{10.1002/1520-6505(2000)9:4<156::aid-evan5>3.0.co;2-7}.

\bibitem[{González-Forero} and {Gardner}(2018)]{gonzalezforero2018}
M.~{González-Forero} and A.~{Gardner}.
\newblock {Inference of ecological and social drivers of human brain-size
  evolution}.
\newblock \emph{Nature}, 557\penalty0 (7706):\penalty0 554--557, 2018.
\newblock \doi{10.1038/s41586-018-0127-x}.

\bibitem[{Dunbar}(1998)]{dunbar1998}
R.~I.~M. {Dunbar}.
\newblock {The social brain hypothesis}.
\newblock \emph{Evol. Anthropol.}, 6\penalty0 (5):\penalty0 178--190, 1998.
\newblock \doi{10.1002/(sici)1520-6505(1998)6:5<178::aid-evan5>3.0.co;2-8}.

\bibitem[{Muthukrishna} et~al.(2018){Muthukrishna}, {Doebeli}, {Chudek}, and
  {Henrich}]{muthukrishna2018}
M.~{Muthukrishna}, M.~{Doebeli}, M.~{Chudek}, and J.~{Henrich}.
\newblock {The Cultural Brain Hypothesis: How culture drives brain expansion,
  sociality, and life history}.
\newblock \emph{PLoS Comput. Biol.}, 14\penalty0 (11):\penalty0 e1006504, 2018.
\newblock \doi{10.1371/journal.pcbi.1006504}.

\bibitem[{Flinn} et~al.(2005){Flinn}, {Geary}, and {Ward}]{flinn2005}
M.~V. {Flinn}, D.~C. {Geary}, and C.~V. {Ward}.
\newblock {Ecological dominance, social competition, and coalitionary arms
  races}.
\newblock \emph{Evol. Hum. Behav.}, 26\penalty0 (1):\penalty0 10--46, 2005.
\newblock \doi{10.1016/j.evolhumbehav.2004.08.005}.

\bibitem[{Miller}(2000{\natexlab{b}})]{miller2000book}
G.~F. {Miller}.
\newblock \emph{{The Mating Mind: How Sexual Choice Shaped the Evolution of
  Human Nature}}.
\newblock Doubleday, New York, 2000{\natexlab{b}}.

\bibitem[{West-Eberhard}(1983)]{westeberhard1983}
M.~J. {West-Eberhard}.
\newblock {Sexual selection, social competition, and speciation}.
\newblock \emph{Q. Rev. Biol.}, 58\penalty0 (2):\penalty0 155--183, 1983.
\newblock \doi{10.1086/413215}.

\bibitem[{Nesse}(2007)]{nesse2007}
R.~M. {Nesse}.
\newblock {Runaway social selection for displays of partner value and
  altruism}.
\newblock \emph{Biol. Theory}, 2\penalty0 (2):\penalty0 143--155, 2007.
\newblock \doi{10.1162/biot.2007.2.2.143}.

\bibitem[{Koster} et~al.(2020){Koster}, {McElreath}, {Hill}, {Yu}, {Shepard},
  {van Vliet}, et~al.]{koster2020}
J.~{Koster}, R.~{McElreath}, K.~{Hill}, D.~{Yu}, G.~{Shepard}, N.~{van Vliet},
  et~al.
\newblock {The life history of human foraging: Cross-cultural and individual
  variation}.
\newblock \emph{Sci. Adv.}, 6\penalty0 (26):\penalty0 eaax9070, 2020.
\newblock \doi{10.1126/sciadv.aax9070}.

\bibitem[{Gurven} et~al.(2006){Gurven}, {Kaplan}, and {Gutierrez}]{gurven2006}
M.~{Gurven}, H.~{Kaplan}, and M.~{Gutierrez}.
\newblock {How long does it take to become a proficient hunter? Implications
  for the evolution of extended development and long life span}.
\newblock \emph{J. Hum. Evol.}, 51\penalty0 (5):\penalty0 454--470, 2006.
\newblock \doi{10.1016/j.jhevol.2006.05.003}.

\bibitem[{Lande}(1980)]{lande1980}
R.~{Lande}.
\newblock {Sexual dimorphism, sexual selection, and adaptation in polygenic
  characters}.
\newblock \emph{Evolution}, 34\penalty0 (2):\penalty0 292--305, 1980.
\newblock \doi{10.1111/j.1558-5646.1980.tb04817.x}.

\bibitem[{Du} et~al.(2018){Du}, {Zipkin}, {Hatala}, {Renner}, {Baker},
  {Bianchi}, {Bernal}, and {Wood}]{du2018}
A.~{Du}, A.~M. {Zipkin}, K.~G. {Hatala}, E.~{Renner}, J.~L. {Baker},
  S.~{Bianchi}, K.~H. {Bernal}, and B.~A. {Wood}.
\newblock {Pattern and process in hominin brain size evolution are
  scale-dependent}.
\newblock \emph{Proc. R. Soc. B}, 285\penalty0 (1873):\penalty0 20172738, 2018.
\newblock \doi{10.1098/rspb.2017.2738}.

\bibitem[{Gingerich}(2022)]{gingerich2022}
P.~D. {Gingerich}.
\newblock {Pattern and rate in the Plio-Pleistocene evolution of modern human
  brain size}.
\newblock \emph{Sci. Rep.}, 12\penalty0 (1):\penalty0 11216, 2022.
\newblock \doi{10.1038/s41598-022-15481-3}.

\bibitem[{Villmoare} and {Grabowski}(2022)]{villmoare2022}
B.~{Villmoare} and M.~{Grabowski}.
\newblock {Did the transition to complex societies in the Holocene drive a
  reduction in brain size? A reassessment of the DeSilva et al. (2021)
  hypothesis}.
\newblock \emph{Front. Ecol. Evol.}, 10:\penalty0 963568, 2022.
\newblock \doi{10.3389/fevo.2022.963568}.

\bibitem[{Lande}(1976)]{lande1976}
R.~{Lande}.
\newblock {Natural selection and random genetic drift in phenotypic evolution}.
\newblock \emph{Evolution}, 30\penalty0 (2):\penalty0 314--334, 1976.
\newblock \doi{10.1111/j.1558-5646.1976.tb00911.x}.

\bibitem[{Li} and {Durbin}(2011)]{lidurbin2011}
H.~{Li} and R.~{Durbin}.
\newblock {Inference of human population history from individual whole-genome sequences}.
\newblock \emph{Nature}, 475\penalty0 (7357):\penalty0 493--496, 2011.
\newblock \doi{10.1038/nature10231}.

\bibitem[{Püschel} et~al.(2024){Püschel}, {Nicholson}, {Baker}, {Barton}, and
  {Venditti}]{puschel2024}
T.~A. {Püschel}, S.~L. {Nicholson}, J.~{Baker}, R.~A. {Barton}, and
  C.~{Venditti}.
\newblock {Hominin brain size increase has emerged from within-species
  encephalization}.
\newblock \emph{Proc. Natl. Acad. Sci. USA}, 121\penalty0 (49):\penalty0
  e2409542121, 2024.
\newblock \doi{10.1073/pnas.2409542121}.

\bibitem[{Hubbe} and {Harvati}(2026)]{hubbe2026}
M.~{Hubbe} and K.~{Harvati}.
\newblock {Evolutionary drivers of encephalization and facial reduction in the
  genus Homo}.
\newblock \emph{Nat. Commun.}, 17\penalty0 (1):\penalty0 5625, 2026.
\newblock \doi{10.1038/s41467-026-74739-w}.

\bibitem[{Berger} et~al.(2015){Berger}, {Hawks}, {de Ruiter}, {Churchill},
  {Schmid}, {Delezene}, et~al.]{berger2015}
L.~R. {Berger}, J.~{Hawks}, D.~J. {de Ruiter}, S.~E. {Churchill}, P.~{Schmid},
  L.~K. {Delezene}, et~al.
\newblock {\textit{Homo naledi}, a new species of the genus \textit{Homo} from
  the Dinaledi Chamber, South Africa}.
\newblock \emph{eLife}, 4:\penalty0 e09560, 2015.
\newblock \doi{10.7554/elife.09560}.

\bibitem[{Dirks} et~al.(2017){Dirks}, {Roberts}, {Hilbert-Wolf}, {Kramers},
  {Hawks}, {Dosseto}, et~al.]{dirks2017}
P.~H. {Dirks}, E.~M. {Roberts}, H.~{Hilbert-Wolf}, J.~D. {Kramers}, J.~{Hawks},
  A.~{Dosseto}, et~al.
\newblock {The age of \textit{Homo naledi} and associated sediments in the
  Rising Star Cave, South Africa}.
\newblock \emph{eLife}, 6:\penalty0 e24231, 2017.
\newblock \doi{10.7554/elife.24231}.

\bibitem[{Sutikna} et~al.(2016){Sutikna}, {Tocheri}, {Morwood}, {Saptomo},
  {Jatmiko}, {Awe}, et~al.]{sutikna2016}
T.~{Sutikna}, M.~W. {Tocheri}, M.~J. {Morwood}, E.~W. {Saptomo}, {Jatmiko},
  R.~D. {Awe}, et~al.
\newblock {Revised stratigraphy and chronology for Homo floresiensis at Liang
  Bua in Indonesia}.
\newblock \emph{Nature}, 532\penalty0 (7599):\penalty0 366--369, 2016.
\newblock \doi{10.1038/nature17179}.

\bibitem[{Pietschnig} et~al.(2015){Pietschnig}, {Penke}, {Wicherts}, {Zeiler},
  and {Voracek}]{pietschnig2015}
J.~{Pietschnig}, L.~{Penke}, J.~M. {Wicherts}, M.~{Zeiler}, and M.~{Voracek}.
\newblock {Meta-analysis of associations between human brain volume and
  intelligence differences: How strong are they and what do they mean?}
\newblock \emph{Neurosci. Biobehav. Rev.}, 57:\penalty0 411--432, 2015.
\newblock \doi{10.1016/j.neubiorev.2015.09.017}.

\bibitem[{Nave} et~al.(2019){Nave}, {Jung}, {Karlsson Linnér}, {Kable}, and
  {Koellinger}]{nave2019}
G.~{Nave}, W.~H. {Jung}, R.~{Karlsson Linnér}, J.~W. {Kable}, and P.~D.
  {Koellinger}.
\newblock {Are bigger brains smarter? Evidence from a large-scale preregistered
  study}.
\newblock \emph{Psychol. Sci.}, 30\penalty0 (1):\penalty0 43--54, 2019.
\newblock \doi{10.1177/0956797618808470}.

\bibitem[{Ponce de León} et~al.(2021){Ponce de León}, {Bienvenu}, {Marom},
  {Engel}, {Tafforeau}, {Alatorre Warren}, et~al.]{poncedeleon2021}
M.~S. {Ponce de León}, T.~{Bienvenu}, A.~{Marom}, S.~{Engel}, P.~{Tafforeau},
  J.~L. {Alatorre Warren}, et~al.
\newblock {The primitive brain of early \textit{Homo}}.
\newblock \emph{Science}, 372\penalty0 (6538):\penalty0 165--171, 2021.
\newblock \doi{10.1126/science.aaz0032}.

\bibitem[{Neubauer} et~al.(2018){Neubauer}, {Hublin}, and {Gunz}]{neubauer2018}
S.~{Neubauer}, J.-J. {Hublin}, and P.~{Gunz}.
\newblock {The evolution of modern human brain shape}.
\newblock \emph{Sci. Adv.}, 4\penalty0 (1):\penalty0 eaao5961, 2018.
\newblock \doi{10.1126/sciadv.aao5961}.

\bibitem[{Deary} et~al.(2000){Deary}, {Whalley}, {Lemmon}, {Crawford}, and {Starr}]{deary2000}
I.~J. {Deary}, L.~J. {Whalley}, H.~{Lemmon}, J.~R. {Crawford}, and J.~M. {Starr}.
\newblock {The stability of individual differences in mental ability from childhood to old age: Follow-up of the 1932
  Scottish Mental Survey}.
\newblock \emph{Intelligence}, 28\penalty0 (1):\penalty0 49--55, 2000.
\newblock \doi{10.1016/S0160-2896(99)00031-8}.

\bibitem[{Borkenau} et~al.(2004){Borkenau}, {Mauer}, {Riemann}, {Spinath}, and
  {Angleitner}]{borkenau2004}
P.~{Borkenau}, N.~{Mauer}, R.~{Riemann}, F.~M. {Spinath}, and A.~{Angleitner}.
\newblock {Thin slices of behavior as cues of personality and intelligence}.
\newblock \emph{J. Pers. Soc. Psychol.}, 86\penalty0 (4):\penalty0 599--614,
  2004.
\newblock \doi{10.1037/0022-3514.86.4.599}.

\bibitem[{Sell} et~al.(2009){Sell}, {Cosmides}, {Tooby}, {Sznycer}, {von
  Rueden}, and {Gurven}]{sell2009}
A.~{Sell}, L.~{Cosmides}, J.~{Tooby}, D.~{Sznycer}, C.~{von Rueden}, and
  M.~{Gurven}.
\newblock {Human adaptations for the visual assessment of strength and fighting
  ability from the body and face}.
\newblock \emph{Proc. R. Soc. B}, 276\penalty0 (1656):\penalty0 575--584, 2009.
\newblock \doi{10.1098/rspb.2008.1177}.

\bibitem[{Connelly} and {Ones}(2010)]{connelly2010}
B.~S. {Connelly} and D.~S. {Ones}.
\newblock {An other perspective on personality: Meta-analytic integration of
  observers' accuracy and predictive validity.}
\newblock \emph{Psychological Bulletin}, 136\penalty0 (6):\penalty0 1092--1122,
  2010.
\newblock \doi{10.1037/a0021212}.

\bibitem[{Smith} et~al.(2017){Smith}, {Schlaepfer}, {Major}, {Dyble}, {Page},
  {Thompson}, et~al.]{smith2017}
D.~{Smith}, P.~{Schlaepfer}, K.~{Major}, M.~{Dyble}, A.~E. {Page},
  J.~{Thompson}, et~al.
\newblock {Cooperation and the evolution of hunter-gatherer storytelling}.
\newblock \emph{Nat. Commun.}, 8\penalty0 (1):\penalty0 1853, 2017.
\newblock \doi{10.1038/s41467-017-02036-8}.

\bibitem[{Spearman}(1910)]{spearman1910}
C.~{Spearman}.
\newblock {Correlation calculated from faulty data}.
\newblock \emph{Br. J. Psychol.}, 3\penalty0 (3):\penalty0 271--295, 1910.
\newblock \doi{10.1111/j.2044-8295.1910.tb00206.x}.

\bibitem[{Brown}(1910)]{brown1910}
W.~{Brown}.
\newblock {Some experimental results in the correlation of mental abilities}.
\newblock \emph{Br. J. Psychol.}, 3\penalty0 (3):\penalty0 296--322, 1910.
\newblock \doi{10.1111/j.2044-8295.1910.tb00207.x}.

\bibitem[{Hill} et~al.(2011){Hill}, {Walker}, {Božičević}, {Eder},
  {Headland}, {Hewlett}, et~al.]{hill2011}
K.~R. {Hill}, R.~S. {Walker}, M.~{Božičević}, J.~{Eder}, T.~{Headland},
  B.~{Hewlett}, et~al.
\newblock {Co-residence patterns in hunter-gatherer societies show unique human
  social structure}.
\newblock \emph{Science}, 331\penalty0 (6022):\penalty0 1286--1289, 2011.
\newblock \doi{10.1126/science.1199071}.

\bibitem[{Driebe} et~al.(2021){Driebe}, {Sidari}, {Dufner}, {von der Heiden},
  {Bürkner}, {Penke}, {Zietsch}, and {Arslan}]{driebe2021}
J.~C. {Driebe}, M.~J. {Sidari}, M.~{Dufner}, J.~M. {von der Heiden}, P.~C.
  {Bürkner}, L.~{Penke}, B.~P. {Zietsch}, and R.~C. {Arslan}.
\newblock {Intelligence can be detected but is not found attractive in videos
  and live interactions}.
\newblock \emph{Evol. Hum. Behav.}, 42\penalty0 (6):\penalty0 507--516, 2021.
\newblock \doi{10.1016/j.evolhumbehav.2021.05.002}.

\bibitem[{Lew-Levy} et~al.(2017){Lew-Levy}, {Reckin}, {Lavi},
  {Cristóbal-Azkarate}, and {Ellis-Davies}]{lewlevy2017}
S.~{Lew-Levy}, R.~{Reckin}, N.~{Lavi}, J.~{Cristóbal-Azkarate}, and
  K.~{Ellis-Davies}.
\newblock {How do hunter-gatherer children learn subsistence skills? A
  meta-ethnographic review}.
\newblock \emph{Hum. Nat.}, 28\penalty0 (4):\penalty0 367--394, 2017.
\newblock \doi{10.1007/s12110-017-9302-2}.

\bibitem[{Deaner} and {Smith}(2013)]{deaner2013}
R.~O. {Deaner} and B.~A. {Smith}.
\newblock {Sex differences in sports across 50 societies}.
\newblock \emph{Cross-Cult. Res.}, 47\penalty0 (3):\penalty0 268--309, 2013.
\newblock \doi{10.1177/1069397112463687}.

\bibitem[{Wiessner}(2014)]{wiessner2014}
P.~W. {Wiessner}.
\newblock {Embers of society: Firelight talk among the Ju/'hoansi Bushmen}.
\newblock \emph{Proc. Natl. Acad. Sci. USA}, 111\penalty0 (39):\penalty0
  14027--14035, 2014.
\newblock \doi{10.1073/pnas.1404212111}.

\bibitem[{Crow}(1958)]{crow1958}
J.~F. {Crow}.
\newblock {Some possibilities for measuring selection intensities in man}.
\newblock \emph{Hum. Biol.}, 30\penalty0 (1):\penalty0 1--13, 1958.

\bibitem[{Price}(1970)]{price1970}
G.~R. {Price}.
\newblock {Selection and covariance}.
\newblock \emph{Nature}, 227\penalty0 (5257):\penalty0 520--521, 1970.
\newblock \doi{10.1038/227520a0}.

\bibitem[{Sackett} et~al.(1988){Sackett}, {Zedeck}, and {Fogli}]{sackett1988}
P.~R. {Sackett}, S.~{Zedeck}, and L.~{Fogli}.
\newblock {Relations between measures of typical and maximum job performance.}
\newblock \emph{Journal of Applied Psychology}, 73\penalty0 (3):\penalty0
  482--486, 1988.
\newblock \doi{10.1037/0021-9010.73.3.482}.

\bibitem[{Beus} and {Whitman}(2012)]{beus2012}
J.~M. {Beus} and D.~S. {Whitman}.
\newblock {The relationship between typical and maximum performance: A
  meta-analytic examination}.
\newblock \emph{Human Performance}, 25\penalty0 (5):\penalty0 355--376, 2012.
\newblock \doi{10.1080/08959285.2012.721831}.

\bibitem[{Reader} et~al.(2011){Reader}, {Hager}, and {Laland}]{reader2011}
S.~M. {Reader}, Y.~{Hager}, and K.~N. {Laland}.
\newblock {The evolution of primate general and cultural intelligence}.
\newblock \emph{Philosophical Transactions of the Royal Society B: Biological
  Sciences}, 366\penalty0 (1567):\penalty0 1017--1027, 2011.
\newblock \doi{10.1098/rstb.2010.0342}.

\bibitem[{Burkart} et~al.(2017){Burkart}, {Schubiger}, and {van
  Schaik}]{burkart2017}
J.~M. {Burkart}, M.~N. {Schubiger}, and C.~P. {van Schaik}.
\newblock {The evolution of general intelligence}.
\newblock \emph{Behav. Brain Sci.}, 40:\penalty0 e195, 2017.
\newblock \doi{10.1017/s0140525x16000959}.

\bibitem[{Zohar} et~al.(2022){Zohar}, {Alperson-Afil}, {Goren-Inbar},
  {Prévost}, {Tütken}, {Sisma-Ventura}, {Hershkovitz}, and
  {Najorka}]{zohar2022}
I.~{Zohar}, N.~{Alperson-Afil}, N.~{Goren-Inbar}, M.~{Prévost}, T.~{Tütken},
  G.~{Sisma-Ventura}, I.~{Hershkovitz}, and J.~{Najorka}.
\newblock {Evidence for the cooking of fish 780,000 years ago at Gesher Benot
  Ya'aqov, Israel}.
\newblock \emph{Nat. Ecol. Evol.}, 6\penalty0 (12):\penalty0 2016--2028, 2022.
\newblock \doi{10.1038/s41559-022-01910-z}.

\bibitem[{Arnold} and {Wade}(1984)]{arnold1984}
S.~J. {Arnold} and M.~J. {Wade}.
\newblock {On the measurement of natural and sexual selection: Theory}.
\newblock \emph{Evolution}, 38\penalty0 (4):\penalty0 709--719, 1984.
\newblock \doi{10.1111/j.1558-5646.1984.tb00344.x}.

\bibitem[{Brown} et~al.(2009){Brown}, {Laland}, and {Borgerhoff
  Mulder}]{brown2009}
G.~R. {Brown}, K.~N. {Laland}, and M.~{Borgerhoff Mulder}.
\newblock {Bateman's principles and human sex roles}.
\newblock \emph{Trends Ecol. Evol.}, 24\penalty0 (6):\penalty0 297--304, 2009.
\newblock \doi{10.1016/j.tree.2009.02.005}.

\bibitem[{Betzig}(2012)]{betzig2012}
L.~{Betzig}.
\newblock {Means, variances, and ranges in reproductive success: comparative
  evidence}.
\newblock \emph{Evol. Hum. Behav.}, 33\penalty0 (4):\penalty0 309--317, 2012.
\newblock \doi{10.1016/j.evolhumbehav.2011.10.008}.

\bibitem[{von Rueden} and {Jaeggi}(2016)]{vonrueden2016}
C.~R. {von Rueden} and A.~V. {Jaeggi}.
\newblock {Men's status and reproductive success in 33 nonindustrial societies:
  Effects of subsistence, marriage system, and reproductive strategy}.
\newblock \emph{Proc. Natl. Acad. Sci. USA}, 113\penalty0 (39):\penalty0
  10824--10829, 2016.
\newblock \doi{10.1073/pnas.1606800113}.

\bibitem[{Sear} and {Mace}(2008)]{sear2008}
R.~{Sear} and R.~{Mace}.
\newblock {Who keeps children alive? A review of the effects of kin on child
  survival}.
\newblock \emph{Evol. Hum. Behav.}, 29\penalty0 (1):\penalty0 1--18, 2008.
\newblock \doi{10.1016/j.evolhumbehav.2007.10.001}.

\bibitem[{Hrdy}(2009)]{hrdy2009}
S.~B. {Hrdy}.
\newblock \emph{{Mothers and Others: The Evolutionary Origins of Mutual Understanding}}.
\newblock Harvard University Press, Cambridge, MA, 2009.

\bibitem[{Boehm}(1999)]{boehm1999}
C.~{Boehm}.
\newblock \emph{{Hierarchy in the Forest: The Evolution of Egalitarian
  Behavior}}.
\newblock Harvard University Press, Cambridge, MA, 1999.

\bibitem[{Marlowe}(2004)]{marlowe2004}
F.~W. {Marlowe}.
\newblock {Mate preferences among Hadza hunter-gatherers}.
\newblock \emph{Hum. Nat.}, 15\penalty0 (4):\penalty0 365--376, 2004.
\newblock \doi{10.1007/s12110-004-1014-8}.

\bibitem[{Smith}(2004)]{smith2004}
E.~A. {Smith}.
\newblock {Why do good hunters have higher reproductive success?}
\newblock \emph{Hum. Nat.}, 15\penalty0 (4):\penalty0 343--364, 2004.
\newblock \doi{10.1007/s12110-004-1013-9}.

\bibitem[{Lidborg} et~al.(2022){Lidborg}, {Cross}, and
  {Boothroyd}]{lidborg2022}
L.~H. {Lidborg}, C.~P. {Cross}, and L.~G. {Boothroyd}.
\newblock {A meta-analysis of the association between male dimorphism and
  fitness outcomes in humans}.
\newblock \emph{eLife}, 11:\penalty0 e65031, 2022.
\newblock \doi{10.7554/elife.65031}.

\bibitem[{Courtiol} et~al.(2012){Courtiol}, {Pettay}, {Jokela}, {Rotkirch}, and
  {Lummaa}]{courtiol2012}
A.~{Courtiol}, J.~E. {Pettay}, M.~{Jokela}, A.~{Rotkirch}, and V.~{Lummaa}.
\newblock {Natural and sexual selection in a monogamous historical human
  population}.
\newblock \emph{Proc. Natl. Acad. Sci. USA}, 109\penalty0 (21):\penalty0
  8044--8049, 2012.
\newblock \doi{10.1073/pnas.1118174109}.

\bibitem[{Karmin} et~al.(2015){Karmin}, {Saag}, {Vicente}, {Sayres}, {Järve},
  {Talas}, et~al.]{karmin2015}
M.~{Karmin}, L.~{Saag}, M.~{Vicente}, M.~A.~W. {Sayres}, M.~{Järve}, U.~G.
  {Talas}, et~al.
\newblock {A recent bottleneck of Y chromosome diversity coincides with a
  global change in culture}.
\newblock \emph{Genome Res.}, 25\penalty0 (4):\penalty0 459--466, 2015.
\newblock \doi{10.1101/gr.186684.114}.

\bibitem[{Gardner} et~al.(2022){Gardner}, {Neville}, {Samocha}, {Barclay},
  {Kolk}, {Niemi}, et~al.]{gardner2022}
E.~J. {Gardner}, M.~D.~C. {Neville}, K.~E. {Samocha}, K.~{Barclay}, M.~{Kolk},
  M.~E.~K. {Niemi}, et~al.
\newblock {Reduced reproductive success is associated with selective constraint
  on human genes}.
\newblock \emph{Nature}, 603\penalty0 (7903):\penalty0 858--863, 2022.
\newblock \doi{10.1038/s41586-022-04549-9}.

\bibitem[{Buunk} et~al.(2008){Buunk}, {Park}, and {Dubbs}]{buunk2008}
A.~P. {Buunk}, J.~H. {Park}, and S.~L. {Dubbs}.
\newblock {Parent--offspring conflict in mate preferences}.
\newblock \emph{Rev. Gen. Psychol.}, 12\penalty0 (1):\penalty0 47--62, 2008.
\newblock \doi{10.1037/1089-2680.12.1.47}.

\bibitem[{Chagnon}(1988)]{chagnon1988}
N.~A. {Chagnon}.
\newblock {Life histories, blood revenge, and warfare in a tribal population}.
\newblock \emph{Science}, 239\penalty0 (4843):\penalty0 985--992, 1988.
\newblock \doi{10.1126/science.239.4843.985}.

\bibitem[{Glowacki} and {Wrangham}(2015)]{glowacki2015}
L.~{Glowacki} and R.~{Wrangham}.
\newblock {Warfare and reproductive success in a tribal population}.
\newblock \emph{Proc. Natl. Acad. Sci. USA}, 112\penalty0 (2):\penalty0
  348--353, 2015.
\newblock \doi{10.1073/pnas.1412287112}.

\bibitem[{Beckerman} et~al.(2009){Beckerman}, {Erickson}, {Yost}, {Regalado},
  {Jaramillo}, {Sparks}, {Iromenga}, and {Long}]{beckerman2009}
S.~{Beckerman}, P.~I. {Erickson}, J.~{Yost}, J.~{Regalado}, L.~{Jaramillo},
  C.~{Sparks}, M.~{Iromenga}, and K.~{Long}.
\newblock {Life histories, blood revenge, and reproductive success among the
  Waorani of Ecuador}.
\newblock \emph{Proc. Natl. Acad. Sci. USA}, 106\penalty0 (20):\penalty0
  8134--8139, 2009.
\newblock \doi{10.1073/pnas.0901431106}.

\bibitem[{Wrangham}(2019)]{wrangham2019}
R.~W. {Wrangham}.
\newblock {Hypotheses for the evolution of reduced reactive aggression in the
  context of human self-domestication}.
\newblock \emph{Front. Psychol.}, 10:\penalty0 1914, 2019.
\newblock \doi{10.3389/fpsyg.2019.01914}.

\bibitem[{Henrich} and {Gil-White}(2001)]{henrich2001}
J.~{Henrich} and F.~J. {Gil-White}.
\newblock {The evolution of prestige: freely conferred deference as a mechanism
  for enhancing the benefits of cultural transmission}.
\newblock \emph{Evol. Hum. Behav.}, 22\penalty0 (3):\penalty0 165--196, 2001.
\newblock \doi{10.1016/s1090-5138(00)00071-4}.

\bibitem[{Hawkes} et~al.(1998){Hawkes}, {O'Connell}, {Jones}, {Alvarez}, and
  {Charnov}]{hawkes1998}
K.~{Hawkes}, J.~F. {O'Connell}, N.~G.~B. {Jones}, H.~{Alvarez}, and E.~L.
  {Charnov}.
\newblock {Grandmothering, menopause, and the evolution of human life
  histories}.
\newblock \emph{Proc. Natl. Acad. Sci. USA}, 95\penalty0 (3):\penalty0
  1336--1339, 1998.
\newblock \doi{10.1073/pnas.95.3.1336}.

\bibitem[{Kramer}(2010)]{kramer2010}
K.~L. {Kramer}.
\newblock {Cooperative breeding and its significance to the demographic success
  of humans}.
\newblock \emph{Annual Review of Anthropology}, 39\penalty0 (1):\penalty0
  417--436, 2010.
\newblock \doi{10.1146/annurev.anthro.012809.105054}.

\bibitem[{McDade} et~al.(2007){McDade}, {Reyes-García}, {Blackinton},
  {Tanner}, {Huanca}, and {Leonard}]{mcdade2007}
T.~W. {McDade}, V.~{Reyes-García}, P.~{Blackinton}, S.~{Tanner}, T.~{Huanca},
  and W.~R. {Leonard}.
\newblock {Ethnobotanical knowledge is associated with indices of child health
  in the Bolivian Amazon}.
\newblock \emph{Proc. Natl. Acad. Sci. USA}, 104\penalty0 (15):\penalty0
  6134--6139, 2007.
\newblock \doi{10.1073/pnas.0609123104}.

\bibitem[{Alami} et~al.(2020){Alami}, {von Rueden}, {Seabright}, {Kraft},
  {Blackwell}, {Stieglitz}, {Kaplan}, and {Gurven}]{alami2020}
S.~{Alami}, C.~{von Rueden}, E.~{Seabright}, T.~S. {Kraft}, A.~D. {Blackwell},
  J.~{Stieglitz}, H.~{Kaplan}, and M.~{Gurven}.
\newblock {Mother's social status is associated with child health in a
  horticulturalist population}.
\newblock \emph{Proc. R. Soc. B}, 287\penalty0 (1922):\penalty0 20192783, 2020.
\newblock \doi{10.1098/rspb.2019.2783}.

\bibitem[{Walter} et~al.(2020){Walter}, {Conroy-Beam}, {Buss}, {Asao},
  {Sorokowska}, {Sorokowski}, et~al.]{walter2020}
K.~V. {Walter}, D.~{Conroy-Beam}, D.~M. {Buss}, K.~{Asao}, A.~{Sorokowska},
  P.~{Sorokowski}, et~al.
\newblock {Sex differences in mate preferences across 45 countries: A
  large-scale replication}.
\newblock \emph{Psychol. Sci.}, 31\penalty0 (4):\penalty0 408--423, 2020.
\newblock \doi{10.1177/0956797620904154}.

\bibitem[{Buss}(1989)]{buss1989}
D.~M. {Buss}.
\newblock {Sex differences in human mate preferences: Evolutionary hypotheses
  tested in 37 cultures}.
\newblock \emph{Behav. Brain Sci.}, 12\penalty0 (1):\penalty0 1--14, 1989.
\newblock \doi{10.1017/s0140525x00023992}.

\bibitem[{Li} et~al.(2002){Li}, {Bailey}, {Kenrick}, and {Linsenmeier}]{li2002}
N.~P. {Li}, J.~M. {Bailey}, D.~T. {Kenrick}, and J.~A.~W. {Linsenmeier}.
\newblock {The necessities and luxuries of mate preferences: Testing the
  tradeoffs.}
\newblock \emph{J. Pers. Soc. Psychol.}, 82\penalty0 (6):\penalty0 947--955,
  2002.
\newblock \doi{10.1037/0022-3514.82.6.947}.

\bibitem[{Murdock}(1967)]{murdock1967}
G.~P. {Murdock}.
\newblock \emph{{Ethnographic Atlas}}.
\newblock University of Pittsburgh Press, Pittsburgh, 1967.

\bibitem[{White}(1986)]{white1986}
D.~R. {White}.
\newblock {Forms and frequencies of polygyny: Standard Sample codes}.
\newblock \emph{World Cultures}, 2\penalty0 (2), 1986.

\bibitem[{Betzig}(1989)]{betzig1989}
L.~{Betzig}.
\newblock {Causes of conjugal dissolution: A cross-cultural study}.
\newblock \emph{Curr. Anthropol.}, 30\penalty0 (5):\penalty0 654--676, 1989.
\newblock \doi{10.1086/203798}.

\bibitem[{Murdock} and {White}(1969)]{murdock1969}
G.~P. {Murdock} and D.~R. {White}.
\newblock {Standard cross-cultural sample}.
\newblock \emph{Ethnology}, 8\penalty0 (4):\penalty0 329--369, 1969.
\newblock \doi{10.2307/3772907}.

\bibitem[{Apostolou}(2007)]{apostolou2007}
M.~{Apostolou}.
\newblock {Sexual selection under parental choice: the role of parents in the
  evolution of human mating}.
\newblock \emph{Evol. Hum. Behav.}, 28\penalty0 (6):\penalty0 403--409, 2007.
\newblock \doi{10.1016/j.evolhumbehav.2007.05.007}.

\bibitem[{Walker} et~al.(2011){Walker}, {Hill}, {Flinn}, and
  {Ellsworth}]{walker2011}
R.~S. {Walker}, K.~R. {Hill}, M.~V. {Flinn}, and R.~M. {Ellsworth}.
\newblock {Evolutionary history of hunter-gatherer marriage practices}.
\newblock \emph{PLoS ONE}, 6\penalty0 (4):\penalty0 e19066, 2011.
\newblock \doi{10.1371/journal.pone.0019066}.

\bibitem[{Whyte}(1978)]{whyte1978}
M.~K. {Whyte}.
\newblock {Cross-cultural codes dealing with the relative status of women}.
\newblock \emph{Ethnology}, 17\penalty0 (2):\penalty0 211--237, 1978.
\newblock \doi{10.2307/3773145}.

\bibitem[{Anderson} et~al.(1992){Anderson}, {Crawford}, {Nadeau}, and {Lindberg}]{anderson1992}
J.~L. {Anderson}, C.~B. {Crawford}, J.~{Nadeau}, and T.~{Lindberg}.
\newblock {Was the Duchess of Windsor right? A cross-cultural review of the socioecology of ideals of female body shape}.
\newblock \emph{Ethol. Sociobiol.}, 13\penalty0 (3):\penalty0 197--227, 1992.
\newblock \doi{10.1016/0162-3095(92)90033-Z}.

\bibitem[{Eastwick} et~al.(2014){Eastwick}, {Luchies}, {Finkel}, and
  {Hunt}]{eastwick2014}
P.~W. {Eastwick}, L.~B. {Luchies}, E.~J. {Finkel}, and L.~L. {Hunt}.
\newblock {The predictive validity of ideal partner preferences: A review and
  meta-analysis.}
\newblock \emph{Psychological Bulletin}, 140\penalty0 (3):\penalty0 623--665,
  2014.
\newblock \doi{10.1037/a0032432}.

\bibitem[{Haselton} and {Buss}(2000)]{haselton2000}
M.~G. {Haselton} and D.~M. {Buss}.
\newblock {Error management theory: A new perspective on biases in cross-sex
  mind reading.}
\newblock \emph{J. Pers. Soc. Psychol.}, 78\penalty0 (1):\penalty0 81--91,
  2000.
\newblock \doi{10.1037/0022-3514.78.1.81}.

\bibitem[{Joel} et~al.(2017){Joel}, {Eastwick}, and {Finkel}]{joel2017}
S.~{Joel}, P.~W. {Eastwick}, and E.~J. {Finkel}.
\newblock {Is romantic desire predictable? Machine learning applied to initial
  romantic attraction}.
\newblock \emph{Psychol. Sci.}, 28\penalty0 (10):\penalty0 1478--1489, 2017.
\newblock \doi{10.1177/0956797617714580}.

\bibitem[{Finkel} and {Eastwick}(2009)]{finkel2009}
E.~J. {Finkel} and P.~W. {Eastwick}.
\newblock {Arbitrary social norms influence sex differences in romantic
  selectivity}.
\newblock \emph{Psychol. Sci.}, 20\penalty0 (10):\penalty0 1290--1295, 2009.
\newblock \doi{10.1111/j.1467-9280.2009.02439.x}.

\bibitem[{Jokela} et~al.(2010){Jokela}, {Rotkirch}, {Rickard}, {Pettay}, and
  {Lummaa}]{jokela2010}
M.~{Jokela}, A.~{Rotkirch}, I.~J. {Rickard}, J.~{Pettay}, and V.~{Lummaa}.
\newblock {Serial monogamy increases reproductive success in men but not in
  women}.
\newblock \emph{Behav. Ecol.}, 21\penalty0 (5):\penalty0 906--912, 2010.
\newblock \doi{10.1093/beheco/arq078}.

\bibitem[{Esteve} et~al.(2016){Esteve}, {Schwartz}, {van Bavel}, {Permanyer},
  {Klesment}, and {García-Román}]{esteve2016}
A.~{Esteve}, C.~R. {Schwartz}, J.~{van Bavel}, I.~{Permanyer}, M.~{Klesment},
  and J.~{García-Román}.
\newblock {The end of hypergamy: Global trends and implications}.
\newblock \emph{Population and Development Review}, 42\penalty0 (4):\penalty0
  615--625, 2016.
\newblock \doi{10.1111/padr.12012}.

\bibitem[{Ueda} et~al.(2020){Ueda}, {Mercer}, {Ghaznavi}, and
  {Herbenick}]{ueda2020}
P.~{Ueda}, C.~H. {Mercer}, C.~{Ghaznavi}, and D.~{Herbenick}.
\newblock {Trends in frequency of sexual activity and number of sexual partners
  among adults aged 18 to 44 years in the us, 2000-2018}.
\newblock \emph{JAMA Network Open}, 3\penalty0 (6):\penalty0 e203833, 2020.
\newblock \doi{10.1001/jamanetworkopen.2020.3833}.

\bibitem[{Rosenfeld} et~al.(2019){Rosenfeld}, {Thomas}, and {Hausen}]{rosenfeld2019}
M.~J. {Rosenfeld}, R.~J. {Thomas}, and S.~{Hausen}.
\newblock {Disintermediating your friends: How online dating in the United States displaces other ways of meeting}.
\newblock \emph{Proc. Natl. Acad. Sci. USA}, 116\penalty0 (36):\penalty0 17753--17758, 2019.
\newblock \doi{10.1073/pnas.1908630116}.

\bibitem[{Bruch} and {Newman}(2018)]{bruch2018}
E.~E. {Bruch} and M.~E.~J. {Newman}.
\newblock {Aspirational pursuit of mates in online dating markets}.
\newblock \emph{Sci. Adv.}, 4\penalty0 (8):\penalty0 eaap9815, 2018.
\newblock \doi{10.1126/sciadv.aap9815}.

\bibitem[{Kleven} et~al.(2019){Kleven}, {Landais}, and {Søgaard}]{kleven2019}
H.~{Kleven}, C.~{Landais}, and J.~E. {Søgaard}.
\newblock {Children and gender inequality: Evidence from Denmark}.
\newblock \emph{American Economic Journal: Applied Economics}, 11\penalty0
  (4):\penalty0 181--209, 2019.
\newblock \doi{10.1257/app.20180010}.

\bibitem[{Hopcroft}(2006)]{hopcroft2006}
R.~L. {Hopcroft}.
\newblock {Sex, status, and reproductive success in the contemporary United
  States}.
\newblock \emph{Evol. Hum. Behav.}, 27\penalty0 (2):\penalty0 104--120, 2006.
\newblock \doi{10.1016/j.evolhumbehav.2005.07.004}.

\bibitem[{Hopcroft}(2021)]{hopcroft2021}
R.~L. {Hopcroft}.
\newblock {High income men have high value as long-term mates in the U.S.:
  personal income and the probability of marriage, divorce, and childbearing in
  the U.S.}
\newblock \emph{Evol. Hum. Behav.}, 42\penalty0 (5):\penalty0 409--417, 2021.
\newblock \doi{10.1016/j.evolhumbehav.2021.03.004}.

\bibitem[{Fieder} and {Huber}(2007)]{fieder2007}
M.~{Fieder} and S.~{Huber}.
\newblock {The effects of sex and childlessness on the association between
  status and reproductive output in modern society}.
\newblock \emph{Evol. Hum. Behav.}, 28\penalty0 (6):\penalty0 392--398, 2007.
\newblock \doi{10.1016/j.evolhumbehav.2007.05.004}.

\bibitem[{Weeden} et~al.(2006){Weeden}, {Abrams}, {Green}, and
  {Sabini}]{weeden2006}
J.~{Weeden}, M.~J. {Abrams}, M.~C. {Green}, and J.~{Sabini}.
\newblock {Do high-status people really have fewer children?: Education,
  income, and fertility in the contemporary U.S.}
\newblock \emph{Hum. Nat.}, 17\penalty0 (4):\penalty0 377--392, 2006.
\newblock \doi{10.1007/s12110-006-1001-3}.

\bibitem[{Kolk}(2023)]{kolk2023}
M.~{Kolk}.
\newblock {The relationship between life-course accumulated income and
  childbearing of Swedish men and women born 1940--70}.
\newblock \emph{Population Studies}, 77\penalty0 (2):\penalty0 197--215, 2023.
\newblock \doi{10.1080/00324728.2022.2134578}.

\bibitem[{Jalovaara} et~al.(2019){Jalovaara}, {Neyer}, {Andersson}, {Dahlberg},
  {Dommermuth}, {Fallesen}, and {Lappegård}]{jalovaara2019}
M.~{Jalovaara}, G.~{Neyer}, G.~{Andersson}, J.~{Dahlberg}, L.~{Dommermuth},
  P.~{Fallesen}, and T.~{Lappegård}.
\newblock {Education, gender, and cohort fertility in the Nordic countries}.
\newblock \emph{European Journal of Population}, 35\penalty0 (3):\penalty0
  563--586, 2019.
\newblock \doi{10.1007/s10680-018-9492-2}.

\bibitem[{Beauchamp}(2016)]{beauchamp2016}
J.~P. {Beauchamp}.
\newblock {Genetic evidence for natural selection in humans in the contemporary
  United States}.
\newblock \emph{Proc. Natl. Acad. Sci. USA}, 113\penalty0 (28):\penalty0
  7774--7779, 2016.
\newblock \doi{10.1073/pnas.1600398113}.

\bibitem[{Kong} et~al.(2017){Kong}, {Frigge}, {Thorleifsson}, {Stefansson},
  {Young}, {Zink}, et~al.]{kong2017}
A.~{Kong}, M.~L. {Frigge}, G.~{Thorleifsson}, H.~{Stefansson}, A.~I. {Young},
  F.~{Zink}, et~al.
\newblock {Selection against variants in the genome associated with educational
  attainment}.
\newblock \emph{Proc. Natl. Acad. Sci. USA}, 114\penalty0 (5):\penalty0
  E727--E732, 2017.
\newblock \doi{10.1073/pnas.1612113114}.

\bibitem[{Hugh-Jones} and {Abdellaoui}(2022)]{hughjones2022}
D.~{Hugh-Jones} and A.~{Abdellaoui}.
\newblock {Human capital mediates natural selection in contemporary humans}.
\newblock \emph{Behav. Genet.}, 52\penalty0 (4-5):\penalty0 205--234, 2022.
\newblock \doi{10.1007/s10519-022-10107-w}.

\bibitem[{Burkart} et~al.(2014){Burkart}, {Allon}, {Amici}, {Fichtel},
  {Finkenwirth}, {Heschl}, et~al.]{burkart2014}
J.~M. {Burkart}, O.~{Allon}, F.~{Amici}, C.~{Fichtel}, C.~{Finkenwirth},
  A.~{Heschl}, et~al.
\newblock {The evolutionary origin of human hyper-cooperation}.
\newblock \emph{Nat. Commun.}, 5\penalty0 (1):\penalty0 4747, 2014.
\newblock \doi{10.1038/ncomms5747}.

\bibitem[{Berna} et~al.(2012){Berna}, {Goldberg}, {Horwitz}, {Brink}, {Holt},
  {Bamford}, and {Chazan}]{berna2012}
F.~{Berna}, P.~{Goldberg}, L.~K. {Horwitz}, J.~{Brink}, S.~{Holt},
  M.~{Bamford}, and M.~{Chazan}.
\newblock {Microstratigraphic evidence of in situ fire in the Acheulean strata
  of Wonderwerk Cave, Northern Cape province, South Africa}.
\newblock \emph{Proc. Natl. Acad. Sci. USA}, 109\penalty0 (20):\penalty0
  E1215--E1220, 2012.
\newblock \doi{10.1073/pnas.1117620109}.

\bibitem[{Dunbar}(2004)]{dunbar2004}
R.~I.~M. {Dunbar}.
\newblock {Gossip in evolutionary perspective}.
\newblock \emph{Rev. Gen. Psychol.}, 8\penalty0 (2):\penalty0 100--110, 2004.
\newblock \doi{10.1037/1089-2680.8.2.100}.

\bibitem[{Crist} et~al.(2016){Crist}, {Dunn-Vaturi}, and {de Voogt}]{crist2016}
W.~{Crist}, A.-E. {Dunn-Vaturi}, and A.~{de Voogt}.
\newblock \emph{{Ancient Egyptians at Play: Board Games Across Borders}}.
\newblock Bloomsbury, London, 2016.

\bibitem[{Akbari} et~al.(2026){Akbari}, {Perry}, {Barton}, {Kariminejad},
  {Gazal}, {Li}, et~al.]{akbari2026}
A.~{Akbari}, A.~{Perry}, A.~R. {Barton}, M.~{Kariminejad}, S.~{Gazal}, Z.~{Li},
  et~al.
\newblock {Ancient DNA reveals pervasive directional selection across West
  Eurasia}.
\newblock \emph{Nature}, 654\penalty0 (8118):\penalty0 419--428, 2026.
\newblock \doi{10.1038/s41586-026-10358-1}.

\bibitem[{Dwork} et~al.(2015){Dwork}, {Feldman}, {Hardt}, {Pitassi}, {Reingold}, and {Roth}]{dwork2015}
C.~{Dwork}, V.~{Feldman}, M.~{Hardt}, T.~{Pitassi}, O.~{Reingold}, and A.~{Roth}.
\newblock {The reusable holdout: Preserving validity in adaptive data analysis}.
\newblock \emph{Science}, 349\penalty0 (6248):\penalty0 636--638, 2015.
\newblock \doi{10.1126/science.aaa9375}.

\bibitem[{Recht} et~al.(2019){Recht}, {Roelofs}, {Schmidt}, and {Shankar}]{recht2019}
B.~{Recht}, R.~{Roelofs}, L.~{Schmidt}, and V.~{Shankar}.
\newblock {Do ImageNet classifiers generalize to ImageNet?}
\newblock In \emph{Proceedings of the 36th International Conference on Machine Learning}, PMLR 97:\penalty0 5389--5400, 2019.

\bibitem[{Campbell}(1979)]{campbell1979}
D.~T. {Campbell}.
\newblock {Assessing the impact of planned social change}.
\newblock \emph{Evaluation and Program Planning}, 2\penalty0 (1):\penalty0
  67--90, 1979.
\newblock \doi{10.1016/0149-7189(79)90048-x}.

\bibitem[{Frayser}(1985)]{frayser1985}
S.~G. {Frayser}.
\newblock \emph{{Varieties of Sexual Experience: An Anthropological Perspective on Human Sexuality}}.
\newblock HRAF Press, New Haven, 1985.

\bibitem[{Healy}(2025)]{healygssr}
K.~{Healy}.
\newblock {gssr: General Social Survey data for use in R}.
\newblock R package, version 0.9. \url{https://kjhealy.github.io/gssr/}, 2025.

\bibitem[{Falconer} and {Mackay}(1996)]{falconer1996}
D.~S. {Falconer} and T.~F.~C. {Mackay}.
\newblock \emph{{Introduction to Quantitative Genetics}}.
\newblock Longman, Harlow, 4th edition, 1996.

\end{thebibliography}
\endgroup

\end{document}
